% Options for packages loaded elsewhere
\PassOptionsToPackage{unicode}{hyperref}
\PassOptionsToPackage{hyphens}{url}
\documentclass[
  12pt,
]{memoir}
\usepackage{xcolor}
\usepackage{amsmath,amssymb}
\setcounter{secnumdepth}{-\maxdimen} % remove section numbering
\usepackage{iftex}
\ifPDFTeX
  \usepackage[T1]{fontenc}
  \usepackage[utf8]{inputenc}
  \usepackage{textcomp} % provide euro and other symbols
\else % if luatex or xetex
  \usepackage{unicode-math} % this also loads fontspec
  \defaultfontfeatures{Scale=MatchLowercase}
  \defaultfontfeatures[\rmfamily]{Ligatures=TeX,Scale=1}
\fi
\usepackage[]{ebgaramond}
\ifPDFTeX\else
  % xetex/luatex font selection
\fi
% Use upquote if available, for straight quotes in verbatim environments
\IfFileExists{upquote.sty}{\usepackage{upquote}}{}
\IfFileExists{microtype.sty}{% use microtype if available
  \usepackage[]{microtype}
  \UseMicrotypeSet[protrusion]{basicmath} % disable protrusion for tt fonts
}{}
\makeatletter
\@ifundefined{KOMAClassName}{% if non-KOMA class
  \IfFileExists{parskip.sty}{%
    \usepackage{parskip}
  }{% else
    \setlength{\parindent}{0pt}
    \setlength{\parskip}{6pt plus 2pt minus 1pt}}
}{% if KOMA class
  \KOMAoptions{parskip=half}}
\makeatother
\usepackage{graphicx}
\makeatletter
\newsavebox\pandoc@box
\newcommand*\pandocbounded[1]{% scales image to fit in text height/width
  \sbox\pandoc@box{#1}%
  \Gscale@div\@tempa{\textheight}{\dimexpr\ht\pandoc@box+\dp\pandoc@box\relax}%
  \Gscale@div\@tempb{\linewidth}{\wd\pandoc@box}%
  \ifdim\@tempb\p@<\@tempa\p@\let\@tempa\@tempb\fi% select the smaller of both
  \ifdim\@tempa\p@<\p@\scalebox{\@tempa}{\usebox\pandoc@box}%
  \else\usebox{\pandoc@box}%
  \fi%
}
% Set default figure placement to htbp
\def\fps@figure{htbp}
\makeatother
\setlength{\emergencystretch}{3em} % prevent overfull lines
\providecommand{\tightlist}{%
  \setlength{\itemsep}{0pt}\setlength{\parskip}{0pt}}
\usepackage{adjustbox}
\usepackage{bookmark}
\IfFileExists{xurl.sty}{\usepackage{xurl}}{} % add URL line breaks if available
\urlstyle{same}
\hypersetup{
  pdftitle={Analytic Introduction to Logic},
  hidelinks,
  pdfcreator={LaTeX via pandoc}}

\title{Analytic Introduction to Logic}
\author{}
\date{}

\begin{document}
\frontmatter
\maketitle

\mainmatter
If you are reading this text, it is because you are taking a course on
\emph{logic}. Most of us understand intuitively that logic has something
to do with thinking correctly, but, because one of the goals of this
course is to improve our capacity for precision and thoroughness in
forming our concepts, we will not be satisfied with a vague and poorly
defined intuition. The question that should immediately arise,
therefore, is ``What precisely is logic?'' This introduction will
attempt to give a precise and rigorous answer to that question; for that
reason, it will need to introduce certain ideas from epistemology,
metaphysics, philosophical anthropology. Although these are not actually
part of subject matter of logic---the \emph{material object}, as we
shall learn to call it---they are nevertheless indispensable for
learning what logic actually \emph{is}. Introducing these topics will
also help serve one of the reasons why logic is taught at the beginning
of philosophical studies at all: namely, to lay the groundwork for the
study of metaphysics.

We will divide this introduction into three chapters: first, an
introduction to the theory behind sciences and arts so as to show that
logic can be classified as both; second an introduction to some basic
concepts of philosophical anthropology, so that can better understand
what an \emph{habitus} (such as a science like logic) is; and finally a
detailed look into the precise definition of logic.

\vfill

\noindent\copyright 2024 Louis Melahn, L.C.\\
This work is licensed under a
\href{http://creativecommons.org/licenses/by-nd/4.0/}{Creative Commons
Attribution-NoDerivatives 4.0 International License}.

\mainmatter

\chapter{Logic as a Science and an
Art}\label{logic-as-a-science-and-an-art}

\section{Sciences vs.~Disciplines}\label{sciences-vs.-disciplines}

Before you read this introduction, chances are that you thought of logic
chiefly as a \emph{discipline}: that is, as a body of knowledge that can
be acquired through study. That is not wrong, exactly; however, it is
limiting. This manner of looking at logic considers it only in the
abstract---that is, it ignores the \emph{being} or the \emph{existence}
of logic as it is actually found in reality. (In the logic of concepts,
we will discuss the difference between \emph{abstract} and
\emph{concrete} concepts.) In fact, this abstract way of looking at a
subject is likely how you have looked at \emph{all} academic fields,
whether that be biology, quantum physics, astronomy, psychology,
economics, or any other course of study. This is a somewhat natural
reduction to make, because we generally gain access to such fields
through formal education in a university setting: that is, by taking
classes, reading books, writing papers, doing laboratory assignments,
and the like.

In reality, however, a field of study does not reduce to an abstract
discipline, a set of propositions that one must acquire. Rather,
acquiring knowledge in a given discipline produces a real change in the
mind---or to use more precise terminology, in the \emph{intellect}---of
the person who learns it. Let me illustrate with an example. Consider
the following equation:
\[i \hbar \frac{d}{d t}\vert\Psi(t)\rangle = \hat H\vert\Psi(t)\rangle\]
Readers familiar with quantum mechanics will recognize this as the
time-dependent form of the Schrödinger equation.\footnote{The equation
  describes the so-called wave function, which is used to determine the
  probability that a given particle is at a given location. Perhaps the
  most familiar use case is for describing the orbitals of electrons
  about an atomic nucleus.} For unfamiliar readers, it is doubtless
completely unintelligible. In short, our human capabilities alone---in
this case our \emph{intellect}, our capacity for knowledge of things in
themselves---are insufficient in order to interpret an equation of this
kind. Rather, in order to understand it, we would need to undergo a
stable change in our intellect. The intellect alone and unaided will not
serve, but it requires a modification or ordering---or, to use the
technical term, an \emph{habitus}.\footnote{I am going to use the Latin
  term \emph{habitus} so as to avoid confusion with the English word
  \emph{habit}, which primarily refers to an involuntary ``tic'' or
  pattern of behavior. An \emph{habitus} (pronounced with a silent
  initial \emph{H}) is a permanent disposition or ordering of something
  to an end.} Since the intellectual \emph{habitus} in question helps us
primarily to know things, rather than to do things, we will call it a
\emph{science}.

Notice that I have not identified the \emph{science} of quantum
mechanics primarily with the \emph{discipline}. Rather, the
\emph{science} of quantum mechanics refers to the permanent and stable
modification of the intellect, the \emph{habitus} that is acquired by
studying quantum mechanics. This is not the usual, modern usage of the
term \emph{science}, but rather, it reflects how classical philosophers
used it---or rather, ἐπιστήμη (\emph{epistéme})\footnote{This term, of
  course, is the root of the word \emph{epistemology}, the study of
  science or knowledge.} for the Greek philosophers, and \emph{scientia}
for Medieval thinkers. We will follow the classical usage of the term
\emph{science} in this introduction: an intellectual \emph{habitus} that
permits or helps us to acquire knowledge. A \emph{scientia} is a real,
existing, concrete attribute of a person---not merely an abstract body
of knowledge, or discipline.

As you can probably tell from the foregoing discussion, logic is also a
science, an ἐπιστήμη: not merely an abstract set of rules, but an actual
intellectual \emph{habitus} that exists (in technical terms, we say that
it is \emph{inherent}) in the logician's very person. It will,
therefore, be useful to explain in greater detail what a science (an
ἐπιστήμη) is, and how it is different from other
\emph{habitus}\footnote{The plural of the Latin word \emph{habitus} is
  also \emph{habitus}.}---and also, how to demonstrate that logic is, in
fact, a science.

\section{The Essential Characteristics of any
Science}\label{the-essential-characteristics-of-any-science}

Taking its cue from Aristotle's seminal work on epistemology, the
\emph{Posterior Analytics}, Medieval thinkers have generally described
science (ἐπιστήμη) as \emph{cognitio certa per causas}:\footnote{See
  \textsc{Aristotle}, \emph{Posterior Analytics} (henceforth cited as
  \emph{PA}), tr. by \textsc{G.R.G. Mure}, I.2, 71b9-11: ``We suppose
  ourselves to possess unqualified scientific knowledge of a thing, as
  opposed to knowing it in the accidental way in which the sophist
  knows, when we think that we know \emph{the cause on which the fact
  depends}, as the cause of that fact and of no other, and, further,
  that the fact could not be other than it is.''} in other words, what
distinguishes a science from mere opinion (Greek: δόξα, \emph{dóxa}) is
that one is able to link the individual facts known to their causes.

For example---one that is perhaps more familiar to readers than quantum
mechanics---consider the example of medical science. For millennia,
humans have been aware that certain diseases are contagious. Perhaps the
most notorious example is during the Black Death in Europe (1347-1351):
in order to avoid the spread of the disease, those who were infected, or
were suspected of being infected, were often
quarantined\footnote{From the French \emph{quarantaine} (``fortyish''),
  which is ultimately derived from the Latin \emph{quadraginta}, forty.}---that
is, isolated from the rest of society for forty days. In other words,
Europeans had formed the correct \emph{opinion} regarding the tendency
of the plague to spread from host to host. However, it was not until the
advent of the microscope and the development of the germ theory of
disease (based largely on the work by Louis Pasteur and Robert Koch in
the late nineteenth century) that the true \emph{causes} of the plague
became known---and thus effective prevention and treatment could be
developed. We could say that the study of infectious disease did not
truly become a \emph{science} until the causes were discovered.
Similarly, even today, a person might be aware that washing one's hands
to avoid the spread of infection is a good idea, or that antibiotics are
an effective treatment against bacterial diseases such as the plague,
and yet not know that infection by bacteria is the \emph{cause}. In that
case, the person has the correct δόξα, or opinion, regarding these
matters, but he does not posses ἐπιστήμη.

Similarly, there is the example of chemistry. For many years, scientists
were aware that certain chemical species could only combine with one in
other in certain proportions. For example, combining one volume of
elemental hydrogen with another volume of elemental chlorine would
produce, once the mixture had cooled to its original temperature, two
volumes of hydrogen chloride (hydrochloric acid). However, combining two
volumes of hydrogen with one volume of chlorine would produce a
three-volume mixture of 1/3 hydrogen and 2/3 hydrogen chloride. The
reason why the hydrogen remained a separate species was a mystery, until
the atomic theory was proposed, and the existence of molecules (stable
combinations of atoms) was discovered. In fact, elemental hydrogen and
chlorine are both constituted by diatomic (two-atom) molecules
(H\textsubscript{2} and Cl\textsubscript{2}), and when they combine,
they form molecules with one hydrogen and one chlorine atom (HCl). This
arrangement explains why some hydrogen remains unreacted: once the
chlorine is entirely consumed, it can no longer react. In other words,
the phenomena observed in the laboratory were not satisfactorily
explained until their \emph{causes} were discovered.

Although the seeking of causes is what distinguishes science from
opinion, Aristotle identifies three other essential characteristics of a
science: namely, it must be \emph{certain}---that is, it must lead to
true knowledge without fail or error\footnote{As Aristotle points out in
  the text cited above, by seeking the causes of a phenomenon, the whole
  purpose is to ensure that ``the fact could not be other than it is''
  (\emph{PA} I.2, 71b11), and as a consequence, such knowledge is free
  from error.}; it must study that which can be attributed to its object
in a \emph{necessary} way (i.e., it identifies what cannot possibly be
otherwise)\footnote{See \emph{PA} I.2, 71b14---15: ``Consequently the
  proper object of unqualified scientific knowledge is something which
  cannot be other than it is.''}; and it must be \emph{universally}
applicable to its object of study.\footnote{See \emph{PA} I.11, 77a5-9:
  ``So demonstration does not necessarily imply the being of Forms nor a
  One beside a Many, but it does necessarily imply the possibility of
  truly predicating one of many; since without this possibility we
  cannot save the universal, and if the universal goes, the middle term
  goes with it, and so demonstration becomes impossible. We conclude,
  then, that there must be a single identical term unequivocally
  predicable of a number of individuals.''}

Again, it can help to think of some examples. Consider the science of
chemistry, which helps us to understand the properties of material
entities, as regards the interactions that occur on the atomic and
molecular level. It is clear that chemistry provides a great degree of
\emph{certainty}: if I mix sodium bicarbonate
(H\textsubscript{2}CO\textsubscript{3}) with hydrochloric acid (a
solution of HCl in water), I can be \emph{sure} that table salt (sodium
chloride, NaCl) and carbon dioxide (CO\textsubscript{2}) will be
produced. Similarly, one way to test to see if I have a base, such as
sodium bicarbonate, is to mix a little hydrochloric acid into it; if it
bubbles, I can be \emph{certain} of it. In a similar vein, when I mix
the acid into the bicarbonate, the outcome is the production of salt and
carbon dioxide---and this result could not possibly be otherwise; that
is, the evolution of these products is a \emph{necessary} result.
Finally, the result is \emph{universal}, because it applies to
\emph{all} similar situations: in short, bicarbonate by nature reacts
with hydrochloric acid, therefore \emph{all} samples of bicarbonate will
react with \emph{all} batches of the acid. Another way of saying this is
that chemistry studies only that which can be attributed \emph{per se}
(in and of itself) or \emph{essentially} to the entities studied.
Finally, as I mentioned above, chemistry is a science precisely because
it seeks to identify the \emph{causes} of these properties and
reactions, which are to be found fundamentally in the atomic and
molecular structures of the materials, especially in the arrangement and
interaction of the electrons about the nuclei of the atoms.

The degree of certainty and necessity---and thus its universal
applicability---can vary from science to science. Clearly, the degree of
certainty and necessity that can be attained in pure mathematics is much
greater than what can be obtained in a social science such as psychology
or economics. Nevertheless, even in cases in which complete certainty,
necessity, or universality are unattainable, a science provides useful
knowledge \emph{to the degree} that it is certain, necessary, and
universal. For example, economics affirms that by and large (but not
necessarily in every single case), when the price of a good or service
rises, the quantity demanded falls. Although economics is unable to
predict what an individual consumer will do, it does accurately account
for general trends among many consumers. What a science cannot do is
account for phenomena that have no essential link to their causes
whatsoever; that is, science cannot account for what is \emph{per
accidens} (completely incidental). To a geneticist studying fruit flies,
for example, what interests him are the laws the govern the inheritance
of genetic characterists---what the flies happen to eat for food does
not interest him, as such, since it does not affect heredity.

In summary: a science or ἐπιστήμη is an intellectual \emph{habitus} that
helps the intellect to understand a given object of study. As such, it
always has four essential characteristics:

\begin{enumerate}
\def\labelenumi{\arabic{enumi}.}
\item
  Most fundamentally, it seeks to connect the phenomena it studies to
  their underlying causes. The seeking of causes is what distinguishes
  science (ἐπιστήμη) from opinion (δόξα), and the other three
  characteristics of a science all follow from this one.
\item
  It provides certain knowledge of its object---that is, it enables the
  one who possesses the science to know without error.
\item
  It accounts for phenomena that cannot be different from the way they
  are, or cannot come about in a different way. That is, it accounts for
  phenomena that are linked necessarily (\emph{per se}), and not merely
  incidentally (\emph{per accidens}), to the object of study.
\item
  Finally, it applies universally to the object of study: there is no
  science, as such, of singulars and particulars.
\end{enumerate}

A couple of things follow from this. First of all, in order to determine
whether a given intellectual \emph{habitus} is a science, it is
sufficient to show that it has these four characteristics. We will
revisit this topic below. Second, it is not possible for science or
ἐπιστήμη, in the strict sense of the term, to be false; that is, false
information is the result, not of science, but of erroneous opinion
(δόξα).

\section{Speculative vs.~Practical
Sciences}\label{speculative-vs.-practical-sciences}

We are not, however quite through with our examination of science: as it
turns out, sciences come in different kinds, depending on their object
of study and method. By now, you may have intuited that by defining
science as the ancient and Medieval philosophers did---as
\emph{scientia} or ἐπιστήμη, that is, as an intellectual \emph{habitus}
whose purpose is to know the properties of a given object of study
according to its causes---we are including under the banner ``science''
many more fields than are customarily placed there. By this definition,
``science'' is not limited to empirical fields such as biology,
chemistry, or physics, but rather includes areas of study as diverse as
mathematics, metaphysics, theology, and ethics. It is sufficient that
each field have a specific and unique object of study (the \emph{genus
subiectum}) and appropriate methods for it.

Sciences can be categorized according to their finality---a criterion
that is particularly useful for us. In the modern parlance, we speak of
\emph{pure sciences} such as biology, chemistry, and physics, and
contrast them to \emph{applied sciences} such as medical science and
engineering. This division can be generalized to apply to all sciences
(in the classical sense, as we have defined the term here): that is,
some sciences seek knowledge of their object for its own sake, whereas
other sciences seek the applications of the knowledge gained from them
to some action. In the classical terminology, the former are called
\emph{speculative sciences}, whereas the latter are called
\emph{practical sciences}.

Again, it can be helpful to look at some examples. Consider the sciences
listed in Table \ref{tab:table1}. What the speculative sciences listed
in the left-hand column have in common is that they all seek knowledge
of the object they study for its own sake. For example, a biologist can
study the behavior of nesting robins without having any practical
purpose whatsoever. This or that scientist may in fact have ulterior
designs (he may be thinking of starting a business raising robins for
their eggs), but that ulterior motive (or lack of one) does not change
the character of his work \emph{as a biologist}. The same can be said
for chemistry, physics, mathematics, and metaphysics\footnote{In fact,
  metaphysics, the study of being (that is, of \emph{ens}, that which
  is) as being, or being \emph{qua} being, is rightly considered the
  most speculative of all sciences, since it has no direct practical
  application.}.

\begin{table}[ht]
\begin{center}
\caption{A few examples of speculative and practical sciences. As much as possible, each practical science is placed next to the speculative science it is derived from.}
\label{tab:table1}
\begin{tabular}{l|l}
\toprule
Speculative & Practical \\
\midrule
Metaphysics & — \\
—           & Ethics \\
Biology     & Medical science \\
Chemistry   & Chemical engineering \\
Physics     & Mechanical engineering \\
Mathematics & — \\
\bottomrule
\end{tabular}
\end{center}
\end{table}

On the other hand, the practical sciences listed in the right-hand
column all have a specific kind of operation (that is, a kind of
action), in mind. What they produce is still a kind of \emph{knowledge},
but in contrast to speculative sciences, for practical sciences,
knowledge is not an end in itself---it is a means to action. For
example, studying physics by itself will not be enough to tell you how
to build a bridge; for that, you would need to study mechanical
engineering. You need to learn how to apply the general principles
discovered in physics to the more concrete problem of how to build
structures. Similarly, studying biology---even human biology---by itself
is insufficient to discover how to cure disease; that is the task of
medical science. As you can probably tell from these examples, many
practical sciences can be pared with a speculative science from which
they are derived. In a similar vein---although it may not be obvious at
first---ethics, the study of human actions (actions that involve the
free choices made by the will) in order to discern whether they are
right or wrong---is also in the column of practical sciences. This is
so, because the entire purpose of ethics is to direct our actions: we
wish to discover which behaviors are compatible with human happiness,
and which ones are not---in other words, the knowledge obtained through
ethics is a means to the end of behaving well.

\section{Is Logic a Science?}\label{is-logic-a-science}

It is now time to turn our attention once more to our
topic---logic---and to apply the principles we have been discussing. As
we saw above, logic is clearly an academic discipline and field of study
(at least when taken in the abstract), but the question now arises: is
logic a \emph{science} like biology, chemistry, and mathematics? In
order to answer this question, we will need to apply our definition: a
science is a \emph{cognitio certa per causas}.

Is logic a \emph{cognitio}---that is, an \emph{habitus}, or stable
disposition, of the intellect? In other words, in the absense of logic,
would the intellect be impeded from acquiring certain types of
knowledge? We will see this more clearly when we give the definition of
logic below, but even based on our intuitive understanding of logic, we
can confidently answer ``yes'': in fact, in the complete absense of
logic, we would be unable to acquire any knowledge \emph{at all}---and a
systematic study of logic makes our thinking easier, more ordered, and
less prone to error.

Does logic provide sure and certain knowledge by linking what it studies
to its causes? As we will see, the answer is again ``yes.'' Logic seeks
the \emph{causes of knowledge}: for example, through logic, I can be
sure that conclusion \(C\) follows from premises \(A\) and \(B\). This
type of causality is different from that sought, say, in the sciences I
listed above, all of which seek the \emph{causes of being}, rather the
causes of knowledge. That is, the \emph{logical} causality proper to
logic is distinct from the \emph{ontological} causality proper to other
sciences---but it is still causality.

Finally, what about the other essential characteristics of a science?
Does logic provide \emph{certain} knowledge---that is, knowledge without
error? Does it supply knowledge that follows necessarily from the nature
of its object of study? Does the knowledge it furnishes apply
universally to its object? Again, the answer to all three questions is
in the affirmative. As regards certainty, the very purpose of logic is
to facilitate our reasoning ability and make it less prone to error. As
regards necessity, again, logic provides the tools we need to see that
the conclusions we reach are the correct ones, and could not be
otherwise. As regards universality, it is clear that logic applies to
\emph{all} of the works of our intellect, without exception.

In fact, we should carefully define our terms, here, because there are
two distinct entities without which knowledge is simply impossible:
there is our intellect, which is the faculty that permits us to make
acts of knowledge---knowledge that captures the very essence and being
of things; and there is the intellectual \emph{habitus} that we are
discussing here. For this reason, intellect, which is sometime called
our \emph{reason} (in Latin: \emph{ratio}; in Greek: λόγος or
\emph{lógos})\footnote{In the strict sense, our \emph{reason} refers to
  our capacity for acquiring new knowledge starting from knowledge that
  we have already acquired: that is, our ability to infer or draw
  conclusions, or, to use the technical term, our \emph{ratiocinium}.
  However, when taken broadly, the term \emph{reason} is synonymous with
  the intellect.}, can be called a kind of ``natural logic.'' However,
logic, in the strict sense of the term, refers to an \emph{habitus}, an
acquired ability that \emph{helps} the intellect in its operation.
However, it should be noted that this logic-as-science, or scientific
logic, can be acquired in an imperfect way (that is, the intuitive grasp
of logic that we acquire from the mere use of our reason), or else in a
``perfect'' or systematic way---which is what this course is attempting
to do.

We must, therefore, conclude that logic is indeed a science.

\section{\texorpdfstring{\emph{Logica utens} vs.~\emph{logica
docens}}{Logica utens vs.~logica docens}}\label{logica-utens-vs.-logica-docens}

Based on our previous discussion, the another question arises: is logic
a \emph{speculative} science or a \emph{practical} science? As we saw,
the purpose of a speculative science is knowledge for its own sake.
Logic certainly seems to fit the bill. St.~Thomas Aquinas in his
commentary on Aristotle's \emph{Posterior Analytics} describes the
purpose of logic in this way:

\begin{quote}
{[}A{]} certain art is required, one that can direct the very act of
reason, through which man, if you will, in that same act of reason may
proceed in an orderly manner, more easily, and without error, and this
art is logic, that is, the rational science.\footnote{\textsc{Thomas
  Aquinas}, \emph{Expositio libri Posteriorum Analyticorum} (henceforth
  cited as \emph{EPA}), proemium 1-2: ``{[}A{]}rs quaedam necessaria
  est, quae sit directiva ipsius actus rationis, per quam scilicet homo
  in ipso actu rationis ordinate, faciliter et sine errore procedat. Et
  haec ars est Logica, idest rationalis scientia'' (my translation).}
\end{quote}

Thus, the purpose of logic is to order the individual actions of our
intellect in such a way that the acquisition of knowledge is easier and
less error-prone. When someone has acquired the science of logic, it
will be much easier for him to distinguish truth from falsehood; in
fact, we could summarize Aquinas' description here as saying that logic
is our instrument for discovering true knowledge, and ensuring that it
is \emph{true}.\footnote{A reminder that, for the ancients,
  \emph{knowledge} (\emph{scientia}, ἐπιστήμη) as such is always true.
  There can be false \emph{opinion} (δόξα), but not false
  \emph{knowledge}; we will follow that convention in this course. (I
  will also mention that the ancients used the same term \emph{scientia}
  or ἐπιστήμη to refer both to the \emph{habitus}---that is, the
  habitual knowledge that is not necessarily in act---and to the
  individual \emph{acts} of knowledge that flow from the
  \emph{habitus}.)} As such, we can categorize logic as a speculative
science, albeit a unique one, because (as we will see) its ultimate goal
is not knowledge of its object---the works of reason---but rather
knowledge of the objects of all the \emph{other} sciences. That is, we
do not study logic in order to study the works of reason for their own
sake; rather, we study the works of reason in order to \emph{use} them
in biology, mathematics, physics, engineering, and every other science
(including logic itself)---indeed in every act of the intellect. We can
summarize this idea by saying that logic is \emph{instrumentally
speculative}: it certainly seeks knowledge for its own sake, but it does
so strictly \emph{as a means or instrument} of obtaining that knowledge.

A new question immediately arises: if logic is a means to obtaining
knowledge, and obtaining knowledge is an action, doesn't that make logic
a practical science? It does indeed. What this means is that---somewhat
surprisingly---``speculative'' and ``practical'' are not mutually
exclusive: a science can be both.

How is this possible? It goes back to what we touched on earlier: a
science, being an intellectual \emph{habitus}, makes it possible for the
intellect to do certain actions that it was unable to do before. An
\emph{habitus} is only required when there is a special difficulty in
acquiring the knowledge in question: for instance (referring to the
example above), I only need to learn (that is, acquire habitual
knowledge of) quantum physics if I need to understand things like the
Schrödinger equation. As we discussed above, many speculative (``pure'')
sciences can be pared with a practical (``applied'') science---for
instance, mechanics (physics) with mechanical engineering, or human
biology with medical science. Why is a distinct applied science
necessary? In short, because the application is not always obvious:
mechanics, for example, tells me how to calculate the forces acting on a
rigid body; it does not tell how to apply this knowledge to building a
bridge. Similarly, human biology tells me that cancers are caused by
genetic mutations that accumulate in the DNA of the body's cells; it
does not tell me the best \emph{treatment} for removing that cancer. In
sum: when the application of speculative knowledge poses a particular
difficulty, a distinct \emph{habitus} (in this case, a distinct
\emph{science}) is required in order to make that application.

With logic, however, the situation is a bit different. Since logic
studies the very works of our reason and (as we will see) the proper
construction of relations between them, merely studying the object of
logic immediately hones our ability to reason. We will return to this
idea below.

In any case, it is clear that, in addition to being a speculative
science, logic is also a practical science: although logic is
speculative insofar as its end goal is knowledge, its manner (or
\emph{mode}) of achieving that goal is practical: said in other terms,
logic teaches us \emph{how to know}. For this reason, logic is said to
be a \emph{modally practical} science.

\section{Is Logic Also an Art
(τέχνη)?}\label{is-logic-also-an-art-ux3c4ux3adux3c7ux3bdux3b7}

Attentive readers will note that Aquinas, when describing the goals of
logic, does a curious thing: he calls it an \emph{art} (\emph{ars}; in
Greek: τέχνη or \emph{téchne}). That is surprising, because we have just
spent the last few pages demonstrating that logic is a \emph{science}.
It may also surprise readers that Aquinas calls it an ``art'' along with
painting, music, sculpture, and theater. In order to understand this, it
is best to take a moment to understand what Aristotle and Aquinas meant
by \emph{art} or τέχνη.\footnote{From the Greek word τέχνη, we derive a
  number of English words, including \emph{technique} and
  \emph{technology}.}

For Aquinas, an art could refer to any intellectual \emph{habitus} that
results immediately in action: in the classical formulation, \emph{recta
ratio factibilium}, right reason applied to things that are made. In
other words, it is a use of reason or the intellect (\emph{ratio}) that
is properly ordered (\emph{recta} or ``right'') to the accomplishment of
determined operations or actions (\emph{factibilium}).

Again, some examples are helpful: suppose you are given a piece of
marble, and instructions to carve it into a statue of David. For the
vast majority of readers, the result would most likely not look anything
like Michelangelo's famous statue of David. As we saw before with the
sciences, in order for the sculptor to be able to set chisel to stone
and produce a work of beauty, he must first acquire an
\emph{habitus}---in this case the art of sculpture. I think it is clear
that the same reasoning applies to a violinist, an opera singer, a
painter, and a photographer. However, for Aquinas, \emph{art} is not
limited to what we would call the ``fine arts''; it applies to any
\emph{habitus} that is necessary for skilled action. Therefore,
surgeons, architects, athletes, lawyers, carpenters, writers, orators,
accountants, and entrepreneurs are all ``artists'' by this definition.

Notice that arts can easily be distinguished by the kind of action they
produce. Some arts, called the \emph{mechanical arts}, produce something
that is external: for example, the art of painting results in physical
artwork such as a fresco or canvass, and carpentry results in woodwork
such as tables and chairs. On the other hand, the \emph{liberal arts}
produce something that is internal to the intellect: for instance, the
art of practicing law does not leave a physical work, but rather
produces intangibles such as legal precedents, decisions, motions, and
verdicts.

As with the sciences, arts have certain characteristics in common. Let
us take, for example, the art of carpentry. A carpenter takes a raw
material---in this case, wood---and fashions it into finished products,
such as tables and cabinets. In order to do so, he applies the
techniques that he has learned: how to cut, trim, and smooth out the
wood, putting the pieces together and gluing them in the right way,
adding the right kind of finish, and so on.

We can generalize this idea to every art: the practitioner of an art
takes some indeterminate entity---its ``raw material''---something that
\emph{can} be fashioned as desired, but \emph{need} not be. He
transforms it using means that are determined and universal, so as to
produce the desired product. That is, we could say that, in general, an
artist's task is to take an indeterminate ``matter'' and remove its
indetermination. For instance, when Michelangelo Buonarotti was given
blocks of marble, the number of possible statues he could make from them
was infinite; it was fully \emph{indeterminate}. However, the more he
worked on them, the closer they became to being fully
\emph{determinate}: say, a Pietà.

In fact, any \emph{habitus} with these two characteristics can be
classified as an \emph{art} or τέχνη. With the mechanical arts it is
easy to see, but the same holds for all of the liberal arts: a defense
lawyer must take all of the facts of the case and his knowledge of the
laws to formulate a cogent defense for his client; an orator must take
the infinite possibilities of a language and formulate a beautiful
speech; a composer must take tones, rhythms, dynamics, and
instrumentation, and form a musical composition.

So, what about logic? Could logic be classified as one of the arts? To
summarize what we saw just now, an \emph{habitus} is an art if it
fulfills the following characteristics:

\begin{enumerate}
\def\labelenumi{\arabic{enumi}.}
\item
  It takes an indeterminate matter that does not infallibly attain its
  fully determined form.
\item
  It transforms that matter into a finished product using determinate
  and universal means.
\end{enumerate}

\noindent Does logic qualify? (1) Regarding the first condition, as we
shall see, the task of logic is to take the works of our reason
(concepts, enunciations, and arguments) and to construct relations
between them according to very strict and determined rules. Even though
our reason is ordered to the truth, there is nothing preventing us from
getting the relations between them incorrect: that is, we can fall into
error and falsehood. Therefore, these works of reason constitute the
\emph{indeterminate matter} of logic. (2) Regarding the second
condition, among all of the possible relations among the works of our
reason, logic constructs---using \emph{universal and determined
rules}---the \emph{correct} relations, the ones that will produce
\emph{true} knowledge. It follows that logic is indeed an art: a liberal
art, evidently, as it does not produce a external product.

\section{How Logic Manages to Be All Three at
Once}\label{how-logic-manages-to-be-all-three-at-once}

Ordinarily, speculative sciences, practical sciences, and arts can be
clearly distinguished from one another, even when their subject matter
is closely related. Continuing the example I mentioned above, human
biology is the speculative science at the root of all medical matters;
medical science is the practical science that applies the knowledge of
human biology to the treatment of disease; and medicine is the art that
applies medical science to concrete acts of treatment. Three distinct
\emph{habitus} are at play here, because neither of the two applications
is trivial. A surgeon, for example, can hardly expect to be sufficiently
trained simply by acquiring the speculative science of human biology. In
fact, not even detailed knowledge of medical science is enough: he can
spend years learning \emph{how} surgeries are done, and yet be
completely unprepared for performing an actual surgery. In short, he
does not become a surgeon without first acquiring the art of
surgery---which in this case implies not only speculative and practical
knowledge, but also a training in the coordination of the hands. In
other words, what renders the three \emph{habitus} distinct from one
another is that there is, if you will, a considerable gap to close
between the speculative science and the practical science, and between
the practical science and the art.

With logic, however, we use our intellect to study works produced by
that same intellect. Therefore, the distance that is present in other
fields is absent in logic. We have already noted above that the
speculative aspect of logic, as a science, cannot be really
distinguished from its practical aspect. It is also clear from our
discussion of arts that the \emph{art} of logic cannot be really
distinguished from the \emph{science} of logic. In other words, formal,
scientific instruction in logic---what is traditionally called
\emph{logica docens}\footnote{From the Latin \emph{doceo}, ``to teach.''}---being
the study of an object that is intrinsic to the very intellect that
studies it, results immediately in an improvement in the exercise of
logic---what is traditionally called \emph{logica utens}\footnote{From
  the Latin \emph{utor}, ``to use or employ.''}. Later on in the course,
you will learn to distinguish and relate different kinds of concepts,
how to relate enunciations to one another, and how to form and analyze
syllogisms. You will be able to use those skills immediately; there is
no need for an \emph{additional} science, nor for a distinct \emph{art},
in order to apply them.

\section{How Logic is Necessary}\label{how-logic-is-necessary}

Before closing this chapter, I would like to say a word about the manner
in which logic is a necessity. As Aquinas says, ``the necessary is what
which cannot not be.''\footnote{\textsc{Thomas Aquinas}, \emph{Summa
  Theologiae}, Leon.~vols. 4--12, Typographia Polyglotta S.C. de
  Propaganda Fide, Rome 1888--1906 (henceforth cited as \emph{S.Th.}), I
  q.~82, a.~1, \emph{responsum} (my translation).} In what way can logic
be said to be necessary? Let us take a moment to unpack this idea
further.

Aquinas argues that there are fundamentally four kinds of
necessity:\footnote{Attentive readers familiar with Aristotle will
  observe that each kind of necessity is closely tied to a corresponding
  species of \emph{cause}: agent, formal, material, and final.}

\begin{quote}
In fact, one way that {[}being necessary{]} pertains to something is by
an intrinsic principle, which can be either material, such as when we
say that it is necessary for everything composed from contraries to be
corrupted, or formal, such as when we say that it is necessary for the
three angles of a triangle to be equal to two right angles. This is
natural or absolute necessity.

Another way that something is unable not to be is by an extrinsic
principle, either an end or an agent. For instance, and end is necessary
when something cannot be obtained without it, or else cannot be obtained
well without it, as food is necessary to life and a horse for a journey.
This is called the necessity of end, which is also sometimes called
utility. There is necessity regarding agents when someone is forced by
an agent, so that he is unable to do the contrary. This is the necessity
of coercion.\footnote{\textsc{Thomas Aquinas}, \emph{Summa Theologiae},
  Leon.~vols. 4--12, Typographia Polyglotta S.C. de Propaganda Fide,
  Rome 1888--1906 (henceforth cited as \emph{S.Th.}), I q.~82, a.~1,
  \emph{responsum} (my translation).}
\end{quote}

Let us look at each of these in turn:

\begin{itemize}
\tightlist
\item
  A \emph{material} necessity has to do with the limitations that matter
  places on material beings: for example, fact that they take up space,
  must exist in a single location, and are corruptible.
\item
  A \emph{formal} necessity (also known as a \emph{natural} or
  \emph{absolute} necessity) stems from the substantial or accidental
  forms of an entity. As you will learn in the course, this kind of
  necessity can be of two kinds: it can stem directly from the entity's
  essence, and thus be part of it's very definition, as when we say that
  it is necessary for man to be rational; or else it can be a property
  that stems from that essence, as when we say that man is capable of
  laughing.
\item
  A \emph{final} necessity---which is probably the kind of necessity you
  are most familiar with---refers to the means required in order to
  obtain an end. It is what we generally refer to when we say that we
  ``need'' something. Aquinas observes that this kind of need can be
  \emph{strict} or \emph{moral}.\footnote{Don't be confused by the term
    ``moral'': it is taken in an etymological sense from the Latin term
    for custom (Latin: \emph{mos}, \emph{moris}). The idea is, people
    are accustomed to using suitable means that are not \emph{strictly}
    necessary.

    As regards ethical norms, they are in fact necessities of end,
    because following them is a necessary condition for happiness.
    Nevertheless, not all of them are \emph{strict} necessities.
    following the basic rules of human decency, such as obeying the Ten
    Commandments, is indeed a \emph{strict} necessity for obtaining
    happiness (that is, \emph{not} following them is incompatible with
    happiness); however, some actions---such as entering religious life
    for those who are called to it---are only necessary for
    \emph{greater} happiness. The latter are ``moral'' necessities in
    the sense we are taking the term.} Food is an \emph{strict}
  necessity for life. Taking a car to go somewhere 10 miles away is not
  a strict necessity, but unless you really enjoy walking, taking a car
  is much better. Thus it is a \emph{moral} necessity.
\end{itemize}

What are we to say about the necessity of logic? First of all, as we
saw, logic is essential for obtaining knowledge. \emph{Natural logic},
which is another name for reason itself, is, according to our analysis,
a \emph{strict} necessity for knowing anything, as should be evident. As
for obtaining systematic, scientific knowledge, it should be noted that
such knowledge can be obtained in an imperfect way (in which we
understand only the easiest conclusions of the science and our dominance
of the material is fuzzy) or in a ``perfect'' way (in which all of the
conditions of possessing a science are satisfied, with the necessary
rigor). In order to get an imperfect grasp of a science, natural logic
is sufficient---although having a scientific logic is much better. In a
word, scientific logic is a \emph{moral} necessity of end for this kind
of knowledge. However, scientific logic is a \emph{strict} necessity of
end for ``perfect'' scientific knowledge.

In sum: we conclude from the foregoing discussion that logic is
simultaneously an instrumentally speculative science, a modally
practical science, and an art, whose object is the works of our reason.
Natural logic is a strict necessity for all knowledge whatsoever. Logic
studied scientifically is a great help for obtaining scientific
knowledge, and absolutely essential for ``perfect'' scientific
knowledge. What remains to discuss in greater detail, therefore, is the
\emph{object} of logic, which we will do after a brief discussion of
some key concepts in philosophical anthropology that are helpful for
understanding what we have just gone over.

\chapter{A Primer on the Metaphysics of
Man}\label{a-primer-on-the-metaphysics-of-man}

In this chapter, we will explore some of the essentials of philosophical
anthropology: the philosophical, or metaphysical, study of
mankind.\footnote{We call it \emph{philosophical} anthropology to
  distinguish it from \emph{cultural} anthropology, which is empirical,
  rather than philosophical, in nature.} We will do this for two
reasons. First of all, although this course is designed primarily to
study logic, it is also intended to be the reader's first exposure to
some of the terminology and ideas that will form the basis of their
philosophical inquiry: for example, \emph{ens}, \emph{substance},
\emph{accident}, \emph{faculty}, \emph{habitus}, and \emph{operation}.
Second, by laying the philosophical bases, it will be much easier to
understand the definition of logic, which is our ultimate goal.

\section{What are metaphysics and
anthropology?}\label{what-are-metaphysics-and-anthropology}

Before beginning in earnest, since this is likely the reader's very
first exposure to the sciences of metaphysics and philosophical
anthropology, a short introduction to what we are talking about is in
order.

\emph{Metaphysics} is the study of anything that exists---of
\emph{entities} (in Latin: \emph{entia}), as we will call them---not
according to this or that particular aspect, but simply \emph{because
they are}. That is, according to Aristotle's famous definition,
metaphysics is the study of existing things (\emph{entia}) \emph{qua}
being.\footnote{As Aristotle states in \emph{Metaphysics}, ``There is a
  science which studies Being \emph{qua} Being, and the properties
  inherent in it in virtue of its own nature. This science is not the
  same as any of the so-called particular sciences, for none of the
  others contemplates Being generally \emph{qua} Being; they divide off
  some portion of it and study the attribute of this portion, as do for
  example the mathematical sciences. But since it is for the first
  principles and the most ultimate causes that we are searching, clearly
  they must belong to something in virtue of its own nature''
  (\textsc{Aristotle}, \emph{Metaphysics} {[}henceforth cited simply as
  \emph{Metaph.}{]}, tr. by \textsc{W.D. Ross}, Oxford, Clarendon Press,
  1924, Γ, I.1, 1003a21-28).

  \par

  In languages, like English, that do not distinguish well between
  \emph{being} in the sense of ``something that exists'' (in Latin:
  \emph{ens}) and \emph{being} as a synonym of ``to be'' (in Latin:
  \emph{esse}), this definition is often translated as ``being
  \emph{qua} being'' or ``being as being.'' However, this rendition is
  rather misleading: metaphysics studies existing \emph{things}
  (\emph{entia}) insofar as they are, not primarily that perfection or
  attribute of existing things that we call ``being'' (\emph{esse}).} We
saw above that sciences study their object with the goal not only of
understanding its properties, but with the hope of tracing those
properties all the way back to their causes. As can be seen, metaphysics
is the most general of sciences: no created entity is excluded from it
whatsoever. Therefore, unlike other sciences, it seeks the properties
that can be attributed to \emph{all} entities without
exception,\footnote{We will touch on some of these, as they regard
  logic, when we go over the \emph{transcendentals} in the part of the
  course on the logic of the concept.} and it seeks the causes on which
\emph{all} things depend. It should be noted that metaphysics studies
\emph{created entities}, and thus God is actually \emph{not} the object
of this science. Rather, He is of interest as both the First and
Ultimate extrinsic Cause of all entities:\footnote{Thus, God is not a
  mere entity, in the strict sense of the term, not a mere thing that
  exists: rather, He is Being Itself (Latin: \emph{Ipsum Esse}).} the
First Efficient and Exemplary Cause, and the Ultimate Final
Cause.\footnote{This pattern---that the cause is not the object of a
  science---is common throughout all sciences. For example, consider
  biology, whose object is living things as regards their bodily life.
  Biology links the life processes of living things to their chemical
  causes: for example, the chemical pathways by which organisms extract
  the energy in the food they eat helps to explain their diets and
  eating habits. Biology, as such, studies the \emph{organisms} and
  their properties, not the chemical causes of those properties;
  however, of course it relies a lot on the findings of chemistry to
  understand those properties in greater depth. In a similar way, the
  task of metaphysics is to study created entities, not God; but knowing
  something about God helps the science considerably in its
  understanding of entities.}

Aristotle called this science various names, including the ``first
Science,'' ``theology'', and ``wisdom''---but not ``metaphysics.'' That
name was given first by the compilers of Aristotle's works located in
Alexandria, Egypt, simply because they placed those treatises
\emph{after} Aristotle's works on physics (in Greek: µετὰ τὰ φύσικα,
\emph{metà tà phýsica}). Despite the trivial origin of the name, by an
accident of the Greek language, it is actually quite appropriate: µετά
can also mean ``beyond,'' and it turns out that metaphysics is the only
philosophical discipline that can transcend \emph{beyond} the physical
(that is, the material) world.

When metaphysical principles and findings are applied to mankind, the
resulting science is \emph{philosophical anthropology}, which could be
defined as the study of man, not merely regarding this or that aspect of
man, but \emph{insofar as he is}, or \emph{insofar as he is man}. Like
other applications of metaphysics, philosophical anthropology seeks to
find the properties that apply universally to its object, man, and links
them to their ultimate intrinsic and extrinsic causes (that is, man's
constitutive metaphysical principles, and God as First and Ultimate
extrinsic Cause.) For our purposes, it will be enough to look at some of
the intrinsic causes of man---that is, some of the metaphysical
principles that constitute him.

\section{\texorpdfstring{Entities (\emph{Entia}), Substances, and
Accidents.}{Entities (Entia), Substances, and Accidents.}}\label{entities-entia-substances-and-accidents.}

The first notion that we will introduce is \emph{ens} (or its plural in
Latin \emph{entia}), which is the term that classical and Medieval
philosophers used to translate Aristotle's τὸ ὄν, the present participle
of the Greek verb εἰμί (``to be''). In English, we will be translating
\emph{ens} as ``entity.'' As we discussed above, an entity is simply
something that is, or (in modern parlance) exists\footnote{In
  metaphysics, there is actually a difference in meaning between
  \emph{existence} (Latin: \emph{exsistentia}; in later Latin
  \emph{existentia}, literally ``a coming forth'') and \emph{being}
  (\emph{esse}). Aristotle uses the terms τὸ ὑπάρχειν (\emph{to
  hypárchein}) and τὸ εἶνια (\emph{to eînai}) to render the same idea.
  \emph{Existence} is the mere fact that something is, whereas
  \emph{being} refers to the ontological perfection that renders
  existence possible. Although in modern parlance, the terms have come
  to be practically interchangeable, in Aristotle and St.~Thomas
  Aquinas, they maintain their distinct meaning. Thus, as in most cases
  I am referring to being as a perfection, I will try to use the verb
  \emph{to be} as much as possible when referring to entities.}: as
Aquinas puts it, \emph{id quod est}, ``that which is.'' Aquinas actually
employs an even more precise description: \emph{id quod habet esse}
(``that which has being''), which emphasizes the fact that the term
``being'' (\emph{esse}) refers to a perfection of the entity in
question---a reality that is constitutive of that entity---not a mere
logical fact.

It should be evident that, outside of God, who is a special case (He is
not an entity exactly, but is best characterized as Being Itself,
\emph{Ipsum Esse}), anything that \emph{is} at all, in any way, is an
entity. That is, there are nothing but entities---anything else, simply
is not (does not exist). But what counts as an \emph{ens}? Let us
suppose that there is an oak tree growing in the back yard. Since the
oak tree clearly \emph{is}, no one can doubt that it is an entity
(\emph{ens}). However, the oak tree does not merely ``be,'' as if being
could occur in isolation; it has \emph{attributes} and properties. For
example, it has a height, a width, a weight, a place, textures, colors,
a time, and a configuration (among many, many others). Are these
attributes themselves entities (\emph{entia})? The answer is not as
obvious, since these attributes are so closely linked to their
substrate, or underlying subject, the oak tree. Nevertheless, we can
resolve the issue the way Aristotle did: since each attribute can be
predicated of the oak tree using the verb \emph{to be}, it is a sign
that we intuitively grasp that these attributes enjoy a certain type of
being. The height of the tree, for example, certainly exists (more
precisely, it \emph{is}), for if there \emph{were} no height, it would
be a very strange tree indeed. On the other hand, the attributes of the
tree hardly enjoy the same kind or degree of being as the tree itself.
The oak tree is independent, it is ``by itself'' and ``on its own''---or
as we say in metaphysical terminology, it \emph{subsists}. In contrast,
the attributes do not exist ``on their own''---it would be absurd to
think of height, color, weight, or any of the examples given above,
without some material body to which they correspond---but rather they
are always ``in'' another entity, or as we say, they \emph{inhere} or
are \emph{inherent} in another entity.

By this reasoning, we have discovered the first and easiest division of
entities: some entities, which we will call \emph{substances} (Greek: αἱ
οὐσίαι) are \emph{subsistent}, that is, they exist by themselves.
Others, which we will call \emph{accidents} (Greek: τὰ
συμβεβηκότα,\footnote{Readers familiar with classical Greek will note
  that Aristotle's term is derived from the preposition σύν (``with''),
  used as a prefix, and the verb βαίνο ``to step or walk''---hence the
  verb συμβαίνω. Although the verb in ancient Greek had the primary
  meaning ``to stand with one's feet together,'' it also meant ``to
  stand with or beside,'' which is the meaning that Aristotle is using
  here. The form τὰ συμβεβηκότα (singular: συμβεβηκός) is the perfect
  participle.} \emph{ta symbebekóta}, literally, ``things that stand
alongside'') are mere attributes or modifications of a substance---as we
say, they are merely \emph{inherent} in a substance. Although substance
and accident can, in a sense, be thought of logically as the different
``types'' of entities, they are best thought of as constituent
principles of a thing: the oak tree has an intrinsic principle that
makes it an oak tree---the substance, or, in this context, more often
referred to as the \emph{essence}---and a set of distinct principles,
also intrinsic, that give the tree its various perfections, such as
height, weight, location, and the like.

At this point, it is appropriate to make a note about the term
\emph{substance} (Latin: \emph{substantia}), which can take on slightly
different, but closely related meanings depending on the context. At the
beginning, when we distinguished substance and accident as different
kinds of \emph{entia}, we were considering substances as fully
constituted wholes: for instance, an oak tree, including all of its
perfections or attributes. Medieval philosophers sometimes used the term
``supposit'' (Latin: \emph{suppositum}, Greek: ὑπόστασις,
\emph{hypóstasis}) to denote this meaning of the term \emph{substance}.
As we just saw, however, \emph{substance} can also be thought of as an
intrinsic constitutive principle that is composed with the accidents;
Western European philosophers have generally used the term
\emph{essence} (Latin: \emph{essentia}) to denote this usage. The two
meanings are so closely related that Aristotle uses the very same word,
οὐσία, for both.

Regardless of how we are using the term \emph{substance}---either as
supposit or as essence---it is always the answer to the question ``What
is it?'' Only substances subsist, and thus, only they constitute a full
answer to this question. As Aristotle clarifies, unlike the accidents,
only substance constitutes what a thing is simply \emph{because} it is
(τὸ τὶ ἦν εἶνια, \emph{to ti ên êinai})---not because of some perfection
or attribute.

In sum: all created realities (which includes everything with which we
have habitual contact) are rightly called \emph{entia} or entities,
because they \emph{are}. Substances are the only entities to
\emph{subsist}, or be ``by themselves,'' and thus they are entities par
excellence. Substances are constituted by an intrinsic principle, called
the \emph{essence}, that gives each substance its definition (that is,
the essence makes that substance what it is), and the essence is
composed with various perfections or attributes called \emph{accidents}.
Entities are always real; anything that lacks being entirely is
\emph{not} properly called an entity. (This idea will become important
below, when we talk about \emph{entia rationis}.)

\section{\texorpdfstring{The Intrinsic Composition of Man: Faculties,
\emph{Habitus}, and
Operations}{The Intrinsic Composition of Man: Faculties, Habitus, and Operations}}\label{the-intrinsic-composition-of-man-faculties-habitus-and-operations}

For the moment, that is enough of an introduction to intrinsic
metaphysical principles to begin to apply them to our subject of
interest, which is man---who is the only creature that makes use of
logic: the sub-human creatures are incapable of it, and angels do not
need it. A human being is, like our oak tree, evidently something that
\emph{subsists}: a substance. And therefore, like any substance, he is
composed of an \emph{essence} (an ontological principle that makes him
who he is) and various \emph{accidents}, which are his perfections. His
essence is itself composed of ontological principles that Aristotle
calls \emph{prime matter} and \emph{substantial form}---in man, we know
these as the \emph{body} and the \emph{soul}---a topic that you will
study at length in cosmology (the philosophy of nature) and metaphysics.

For the purposes of our discussion about logic, it is sufficient to
examine certain types of perfections: faculties, \emph{habitus}, and
operations.

\subsection{Faculties}\label{faculties}

It should be obvious from experience that the substances we encounter do
not merely ``be''; they also ``do''---that is, they perform a whole
series of actions (which we will learn to call \emph{operations}). Even
inanimate objects at least exert forces (for instance, the force that a
rock exerts when I stub my toe against it), emit or reflect radiation
(which is how we can see things), and can be made to interact with other
substances (for instance, think of all the chemical reactions that occur
constantly). Evidently, as we go up the scale to nobler creatures, such
as plants, animals, and humans, the number and perfection of the
operations performed increases. For example, consider those substances
that enjoy vegetative life:\footnote{In philosophy, we commonly refer to
  such substances as \emph{plants}, but I would like to emphasize that
  in this context we understand the term in a philosophical, not a
  biological, sense; that is, corporeal substances that are capable of
  growth, nutrition, and reproduction, but do not enjoy sense or
  rational operations. Thus, these substances include many that are not
  in the plant kingdom as biology defines it, such as bacteria, archaea,
  fungi, and protists.} they have available to them a whole series of
operations that are unknown to inanimate objects, such as growth,
nutrition, and reproduction. Animals go a step further: they display
sense knowledge and appetite, and most animals can move from place to
place on their own power. Finally, humans, aptly called \emph{rational
animals}, can do everything that plants and animals can do, but in
addition, we are able to know things in themselves, and to love them for
their own sake.

It may be less obvious at first glance, but, save in the case of
God\footnote{In Thomistic and Aristotelian philosophy, God is perfectly
  identical with both His essence and His operations. God has no
  accidents, or more precisely, every perfection that for a creature
  would be an accident, in God, is His very Essence.}, the raw
essence\footnote{Remember, when we distinguish the \emph{essence} from
  the \emph{substance} in metaphysics, by ``essence'' we mean that
  ontological principle that serves as the subject or substrate of the
  accidents. (Although we have not yet discussed it here, is also the
  potential principle that receives the act of being. You will learn
  more about that in your course on metaphysics.)} of a substance is
incapable of acting or doing anything on its own. In order for the
substance to act in any way, it first needs a \emph{capacity to act} in
that way, and that capacity is itself an accident or perfection of the
substance. For example, in order for an animal to move from place to
place, it must first have a capacity for locomotion (which in higher
animals entails a muscular system and a system of motor nerves).
Similarly, our capacity to know things in themselves, which, as we saw,
is called the \emph{intellect}, is absolutely necessary in order for us
to make acts of knowledge. In fact, all of our operations have their
root in some intrinsic capacity: our physical movement, our senses, our
memory, our imagination, our emotions, our intellectual knowledge, and
our acts of love. In general, any capacity that produces an operation
can be called an \emph{active potency} or \emph{active power}. If the
active potency produces an operation that only an animal (rational or
irrational) can do, we call it a \emph{faculty}.

\begin{table}[ht]
\begin{center}
\caption{Some operations and their corresponding faculties.}
\label{tab:operations-and-faculties}
\begin{adjustbox}{max width=\textwidth}
\begin{tabular}{l|l}
\toprule
Operation & Faculty \\
\midrule
Movement from one place to another & Locomotion \\
Sensation of material entities as regards color & Sight \\
Sensation of material entities as regards sound & Hearing \\
Recollection of past benefits and dangers & Sense memory \\
Composing coherent spatial representations & Common sense \\
Perceiving spatial relations, continuity, and benefit or harm & \emph{Vis cogitativa} \\
Desiring the easy and pleasurable good & Concupiscible appetite \\
Desiring the difficult and arduous good & Irascible appetite \\
Knowing things as they are & Intellect \\
Loving things for their own sake & Will \\
\bottomrule
\end{tabular}
\end{adjustbox}
\end{center}
\end{table}

Since this is a course on logic, it should come as no surprise that the
faculty we are most interested in as the ``cause'' of logic is the
intellect: our capacity to know things as they are. It is worth taking a
moment to understand fully what we are taking about. We should note,
first of all, that it is possible too \emph{know} in different ways. At
its most fundamental level, to \emph{know} something means that the
knower, in a certain sense, or a certain part of him, actually
\emph{becomes} the thing known. We wouldn't normally think of it as a
kind of ``knowledge,'' but in a way the block of marble that
Michelangelo formed into a statue of David came to ``know'' the model
that was in Michelangelo's mind once the chiseling was complete. In a
similar way, when a dog recognizes its owner, it does so by means of the
\emph{sensible image} that it has formed in its internal
senses.\footnote{We now know, through modern biology, that the internal
  senses are located in the central nervous system, especially the
  brain.} This image, in which the accidental forms of the owner have
become impressed on the sense faculties of the dog, can be characterized
in a much more proper sense as \emph{knowledge}: we would call it
\emph{sense knowledge}, the kind of knowledge that can be acquired only
through the internal and external senses, without the need for a
properly rational faculty.

However, in logic we are not directly interested in sense knowledge.
Through sense knowledge, an animal can only gain access to the object of
its knowledge insofar as it corresponds to (or goes against) its
appetites: in short, an animal knows something only \emph{insofar as it
is perceived as beneficial or noxious}. On the other hand, human beings
are able to do something considerably greater: in addition to perceiving
an object as beneficial or noxious, we can grasp \emph{what} the object
is \emph{in itself}. A dog will wag its tail when its owner brings it
food; it is idle to ask the dog \emph{what food is} or \emph{what an
owner is}, still less whether these entities \emph{exist in reality} or
not, but a human being has no trouble doing either.\footnote{C.S. Lewis
  in his book \emph{That Hideous Strength} provides perhaps the best
  description of the difference between sense and rational knowledge
  that I have come across: ``Mr.~Bultitude's {[}a tame bear that figures
  prominently in the story{]} mind was as furry and as unhuman in shape
  as his body. He did not remember, as a man in his situation would have
  remembered, the provincial zoo from which he had escaped during a
  fire, nor his first snarling and terrified arrival at the Manor, nor
  the slow stages whereby he had learned to love and trust its
  inhabitants. He did not know that he loved and trusted them now. He
  did not know that they were people, nor that he was a bear. Indeed he
  did not know that he existed at all: everything that is represented by
  the words \emph{I} and \emph{Me} and \emph{Thou} was absent from his
  mind. When Mrs.~Maggs gave him a tin of golden syrup, as she did every
  Sunday morning, he did not recognise either a giver or a recipient.
  Goodness occurred and he tasted it. And that was all. Hence his loves
  might, if you wished, be all described as cupboard loves: food and
  warmth, hands that caressed, voices that reassured, were their
  objects. But if by a cupboard love you meant something cold or
  calculating you would be quite misunderstanding the real quality of
  the beast's sensations. He was no more like a human egoist than he was
  like a human altruist. There was no prose in his life. The appetencies
  which a human mind might disdain as cupboard loves were for him
  quivering and ecstatic aspirations which absorbed his whole being,
  infinite yearnings, stabbed with the threat of tragedy and shot
  through with the colours of Paradise. One of our race, if plunged back
  for a moment in the warm, trembling, iridescent pool of that
  pre-Adamite consciousness, would have emerged believing that he had
  grasped the absolute: for the states below reason and the states above
  it have, by their common contrast to the life we know, a certain
  superficial resemblance. Sometimes there returns to us from infancy
  the memory of a nameless delight or terror, unattached to any
  delightful or dreadful thing, a potent adjective floating in a
  nounless void, a pure quality. At such moments we have experience of
  the shallows of that pool. But fathoms deeper than any memory can take
  us, right down in the central warmth and dimness, the bear lived all
  its life'' (taken from \textsc{C.S. Lewis}, \emph{That Hideous
  strength}, New York, Scribner 1999, chapter 14, ``Real Life is
  Meeting,'' part III, p.~306).} As a \emph{rational} animal, a human
being is capable of grasping two important aspects of the object of his
knowledge that are unavailable to an irrational animal: its
\emph{essence} (\emph{what} that object is) and its \emph{being}
(\emph{whether} or not, in fact, that object is or exists). As we will
learn, these aspects of rational knowledge form the basis for logic's
object of study, which is the works of our reason. Through our concepts,
we apprehend \emph{essences}, and through our enunciations (and through
the reasoning that we sometimes need to form enunciations) we make
judgments regarding \emph{being}.

The distinction between sense and intellectual knowledge cannot be
overemphasized: the distinction is \emph{not} made in a large number of
philosophers, especially those who derive their ideas, directly or
indirectly, from the Empiricist school. It is not my purpose to
denigrate the valid insights brought forward by great philosophers such
as John Locke, David Hume, and Immanuel Kant; however, I think each of
these authors \emph{can} be constructively critiqued for overlooking
this important point. As you will learn in philosophical anthropology,
our capacity for intellectual knowledge and love (our intellects and
wills), as distinct from sense knowledge and sense appetite, is the most
important evidence in favor of the \emph{immateriality} of our souls.
Our intellects do not depend for their existence on our bodies, but
rather emanate directly from our rational, spiritual, immaterial souls.
What confused the Empiricists, in my opinion, is that our intellects,
emanating from a soul that informs a body, requires sensation as a
necessary condition for its operation (barring a divine intervention).
It is true, as the Scholastic saying aptly states, that \emph{nihil est
in intellectu quod prius non fuerit in sensu} (``there is nothing in the
intellect that was not first in the senses''); nevertheless,
intellectual knowledge is not \emph{produced} by the body (as sense
knowledge is), but rather by a spiritual, immaterial faculty that
transcends the body. Said in other terms, the brain is not the ``organ''
that ``holds'' the intellect; rather, the intellect makes use of the
brain (and the rest of the body's sensory apparatus) in order to
function properly.

This is a good moment to make some more remarks about terminology. The
most correct and proper name for the faculty through which we know
things as they are is the \emph{intellect} (Latin: \emph{intellectus};
Greek: νοῦς, \emph{noûs}). The term ``mind'' (Latin: \emph{mens}) is
less precise, because it can refer also to emotional states. Similarly,
``consciousness'' (a favorite of many contemporary philosophers) tends
to imply that the person in question is actually \emph{exercising} his
intellect---i.e., that he is actively aware of his
surroundings---whereas the intellect exists even when the person is
asleep or otherwise unconscious. (Indeed, it exists even before he has
developed the use of his reason.) Moreover, ``consciousness'' and
``self-awareness'' are difficult to define, and could easily be
attributed to some of the higher animals. Finally, you will often
encounter the term ``reason'' (Latin: \emph{ratio}; Greek: λόγος,
\emph{lógos}) used as a rough synonym for the intellect. When speaking
in a broad sense, it is fine, but strictly speaking the term ``reason''
applies to our ability to take known premises and derive previously
unknown conclusions from them: what in Latin is more precisely called
the \emph{ratiocinium}.\footnote{I will alert the reader from the outset
  that there is an ambiguity in the way that Aquinas uses the Latin term
  \emph{intellectus}. In most cases, he means the faculty by which we
  know things as they are, but in a few cases he is referring to an
  \emph{habitus} by which we know the first principles of reason. We
  will discuss that later on when we talk about the object of logic.} It
should also be noted that only a \emph{human} intellect can properly be
called ``reason'', because the term refers to the fact that the human
intellect is inherently \emph{discursive}: that is, it uses internal
\emph{words} (Greek: \emph{lógoi}; Latin: \emph{rationes} or
\emph{verbi}). The irrational animals do not have any intellect at all,
and the angels (and God) have a strictly intuitive intellect that does
not require discourse.

Now, as I will emphasize many times, logic studies neither the human
intellect nor its operations (both of which would pertain to rational
psychology or philosophical anthropology), but rather the \emph{works
produced} by that intellect, which, as we will see below are of three
kinds: concepts, enunciations, and arguments. Nevertheless, since we are
attempting to study logic in a scientific way, and the intellect is the
\emph{cause} of those works, it is important to understand what the
intellect is and what it does.

In sum: faculties (as part of the larger genus of qualities that we call
active potencies or active powers) are the first accidents to emanate
from a substance. They provide the substance with the basic capacities
to \emph{do} things and in general to \emph{act}. They are the
\emph{sine qua non} for any operation whatsoever. Since we are studying
logic, which studies the works of the human intellect, naturally, the
faculty that we are going to discuss the most is the intellect. We must
now turn to the ways that faculties can, as it were, be ``trained'' (or
deformed) in order to function better (or worse): that is, the moment
has arrived to discuss \emph{habitus} systematically.

\subsection{\texorpdfstring{\emph{Habitus}}{Habitus}}\label{habitus}

We are now in a position to talk about \emph{habitus} more
systematically, having touched on them in the previous chapter. In
general, an \emph{habitus} (Greek: ἕξις, \emph{héxis}) can be thought of
as a stable ordering of a substance, which either perfects it or deforms
it.

Here is an analogy: suppose that I take a clean sheet of white paper. In
itself, the paper has a large number of possible configurations: it
could be flat, rolled into a tube, folded, or crumpled into a ball.
Let's suppose I roll it into a tube, gently, without creasing it. If I
wish to go back to a flat sheet, there is no particular obstacle: all I
need to do is uncurl it.

Now, suppose that instead of rolling it, I \emph{fold} the paper neatly
into three parts, making two creases (the way one folds a letter). Now,
if I ever wanted to go back to a flat sheet, I would have trouble, since
the creases are hard to remove. On the other hand, the paper will never
fit into a letter-sized envelope unless I first fold it.

This analogy helps us to understand what in philosophy we call
\emph{dispositions} (Latin: \emph{dispositiones}; Greek: διαθέσεις,
\emph{diathéseis}) and \emph{habitus}. When I roll up the paper gently,
I do change something about it---I ``perfect'' the paper, as it
were---but the tube-shaped configuration of the paper is transitory.
Similarly, dispositions are transitory perfections of a substance that
can easily be ``undone.'' On the other hand, the creased fold causes a
\emph{permanent} change in the paper that is difficult to undo.
Similarly, an \emph{habitus} creates state in the substance that is
difficult to remove.

Some examples are in order. Returning to the example we made earlier,
all human beings, thanks to their rational nature, have the raw capacity
to understand a science like quantum mechanics: that is, they all have
the faculty that we call the intellect. However, as demonstrated by the
fact that most people do not immediately understand the Schrödinger
equation, our intellect needs ``training'' or formation in order to make
accurate assertions about quantum physics. With some basic formation, we
can acquire the correct \emph{opinion} (\emph{dóxa}) about certain
facts---for instance, perhaps we have heard the idea that in the very
microscopic scale, things behave both as waves and as particles---but
this cannot be said to be a permanent and thorough assimilation. Armed
with just that opinion, we will hardly be able to understand the
two-slit experiment or the Pauli exclusion principle. In order to do
that, one must be able (as we have seen) to trace the phenomena to their
causes---which, in the case of quantum physics would entail deriving the
relevant equations from the data and models that give rise to them. Only
at that point is the assimilation secure and permanent---that is, truly
a \emph{scientia} (\emph{episteme}), as we have understood the term.
Right \emph{opinion} about quantum physics is a disposition that enables
the intellect to make \emph{some} right assertions about the quantum
world, in a transient and often unreliable way; acquiring the
\emph{science} of quantum mechanics, including tracing the phenomena to
their causes, give secure, necessary, and permanent knowledge. Thus, a
science is best described as an \emph{intellectual habitus} that enables
the one who has it to be secure in his knowledge of an object.

Notice that the intellect could also be malformed: I could gain
\emph{false} information about quantum physics, which could be
transitory or deeply engrained. If the latter, it would consist in a
kind of ignorance or prejudice: a deformation of the intellect. If this
defect is durable and difficult to remove, it, too, is an intellectual
\emph{habitus}, but a bad one.

Thus, I have described four ``perfections'' (two of them better
described as \emph{im}perfections) of rational substances that affect
their ability to use their intellect: two dispositions (true opinion,
false opinion) and two \emph{habitus} (science and deep-seated
prejudice) that ``train'' (or malform) the intellect, enabling them to
perform operations that they were unable to do before. When an
\emph{habitus} affects operation like this, we call it an
\emph{operative habitus} (\emph{habitus operativus}); when that
operative \emph{habitus} is a good one (when it actually makes the
substance more perfect) we call it a \emph{virtue} (Latin:
\emph{virtus}, Greek: ἀρετή, \emph{areté}); when it \emph{deforms} that
substance, rendering it less perfect, we call it a \emph{vice} (Latin:
\emph{vitium}).\footnote{There is no special name for good or bad
  dispositions.}

Perhaps you are more familiar with applying the terms \emph{virtue} and
\emph{vice} to morality. We often say that a person is \emph{virtuous}
if he behaves in a way that is in conformity with the natural law. It
would, however, be more precise to speak more specifically: a
\emph{moral virtue} is a good operative \emph{habitus} that stably and
durably helps a person to behave in a morally acceptable way, making
such behavior easier to accomplish and even pleasurable. In contrast, a
\emph{moral vice} is a bad operative \emph{habitus} that stably and
durably pushes the person to commit immoral acts. Each of these types of
\emph{habitus} also has a corresponding \emph{disposition}, which helps
the person to behave in a morally good (or bad) manner---but in a
transient and non-durable way.

Let's look at some concrete examples. Suppose that one day when I am at
work, I find a \$20 bill on my desk. I strongly suspect that it belongs
to my co-worker. I could really use that money, but---after hemming and
hawing for a while---I decide that really the morally good thing to do
would be to return it to my co-worker. But if I am truthful, the
decision was a close one; I nearly pocketed the money. You could say I
performed a morally good act, and act of \emph{honesty}, which is a
species of \emph{justice} (that is, of giving each one his due).
However, I did so reluctantly. Now, suppose that my colleague has a
tendency to lose money, often leaving it carelessly on my desk. Thus, at
least once a day I have to decide whether or not to restore the money or
pocket it. As it turns out, after a week, I have returned the money
every time; and after a month, it becomes practically second nature to
me. Soon, I actually experience pleasure in doing so, especially given
my colleague's gratitude. At the end of this process, we can say that I
have acquired the moral virtue of \emph{honesty}: that is, I now have a
stable and durable tendency to return lost money to its owner.

However, of course, my free human acts are contingent: if I wanted to, I
could act \emph{contrary} to my conscience. Let's go back to the first
day, when I first found a \$20 bill on by desk: suppose that, in this
new timeline, my greed gets the better of me and I decide to pocket the
money. I don't get caught, and my colleague never discovers where he has
left the money. At this point, I have committed a morally unacceptable
act: an act of \emph{dis}honesty (indeed, a type of theft). However, I
have not yet acquired a stable tendency to act this way. Unfortunately,
over the course of time, I take advantage of my colleague's
carelessness, pocketing hundreds of dollars over the course of two
months. I get better at figuring out where he leaves the money, even
scheming to get him to drop \emph{more} money, and I take pleasure in my
newfound ``wealth.'' By now, I have acquired a stable disposition to
theft, thus forming a \emph{vice}.

Note that, in either case---with both the virtue and the vice---there is
a period of transition during which it is still relatively easy for me
to change course. In that intermediate stage, we can say that I have a
\emph{disposition} toward honesty (or dishonesty), but not yet a virtue
(or a vice). Nevertheless, that disposition becomes stronger and more
deeply seated as I continue to favor it with my action, until it
``hardens'' into an \emph{habitus}.

These examples, I think, are sufficient to illustrate how operative
\emph{habitus} (and dispositions) work: in short, they are principles of
operation that have the function of perfecting a faculty, giving it the
ability to do things that it was unable to do before, or else to do them
with greater ease. We can say that they are, as it were, halfway between
the faculty and its operation: they partially actualize or determine the
faculty, causing it to tend toward some operations over others.

Attentive readers will observe that up to now I have consistently
specified that the \emph{habitus} and dispositions we have talked about
so far are all \emph{operative}: that is, they are the principles of
various kinds of \emph{operation}. I have concentrated on the operative
kind, because, of course, logic---our unique intellectual virtue that is
simultaneously a speculative science, a practical science, and an
art---is a intellectual operative \emph{habitus}: evidently, it is a
\emph{good} one, hence, an intellectual virtue. However, I will mention
in passing that there is also a whole class of dispositions and
\emph{habitus} that are not immediately ordered to operation. Rather,
they produce a state of \emph{being} in the substance. These
\emph{habitus} and dispositions are called \emph{entitative}. To give a
few examples, health and physical beauty are both entitative
dispositions; and the sanctifying grace that Christians receive at
Baptism is perhaps the best example of an entitative \emph{habitus}.

\begin{figure}
\centering
\pandocbounded{\includegraphics[keepaspectratio]{images/habitus-examples-en.pdf}}
\caption{The divisions of \emph{habitus} and some examples of each. Note
that dispositions are divided the same way as \emph{habitus}, and, in
fact, physical beauty and health are best described as dispositions.
\label{fig:habitus-examples}}
\end{figure}

It is worth noting at this point some of the elements that operative
\emph{habitus}\footnote{From now on, for brevity's sake I will speak
  only of \emph{habitus}, because that is the genus that logic falls
  under. However, all of these characteristics apply equally well to
  dispositions, unless otherwise specified. As we saw, dispositions are
  like \emph{habitus} in all respects, except for durability.} all have
in common, many of which they share in common with faculties. In
particular, both operative \emph{habitus} and faculties are inherent in
a \emph{subject} or substrate, and both are principles or originators of
\emph{operation}, and both have an \emph{object}, that is, the finality
or term to which the operation is directed.

As usual, it is helpful to look at examples. Let us look first of all at
some faculties, starting with the faculty of sight. Like all faculties,
the \emph{subject} (or \emph{substrate}) will be the substance in which
it is inherent (in this case, an animal, rational or irrational). We can
specify, however, that unlike purely immaterial faculties like the
intellect, a sense like sight is dependent on the \emph{body} or
\emph{matter} of the animal in question. The \emph{operation} that sight
produces is called \emph{vision}, and the \emph{object} of sight can be
described as \emph{colored material bodies insofar as they have color}.
Now, let us look at the human intellect. Being also a faculty, it will
likewise have the substance as its \emph{subject} (in this case, the
human being in question)---however, unlike the sense of sight, it stems
directly from the human being's immaterial soul, with no dependence on
the body for its existence. Its \emph{operations} include the simple
apprehension of essences, and making judgments (mediated or immediate)
regarding being. (We will return to this point in greater detail.) Its
\emph{object} can be described as the entities it comes to know, insofar
as they are intelligible (that is, insofar as they can be known in
themselves).

With operative \emph{habitus}, it is important to specify what we mean
by \emph{subject}: like all accidents, operative \emph{habitus} inhere
in the substance that they perfect, and so the substance could be
regarded as their ``subject'' in that sense. However, more often, when
we speak of the \emph{subject} of an operative \emph{habitus}, we mean
the faculty that the \emph{habitus} modifies. That is the meaning that
we will give the term in this introduction. Thus, for example, all the
intellectual \emph{habitus} that we have talked about up to
now---sciences and arts---have the intellect as their subject. If we
look at, say, quantum mechanics, its \emph{object} is what the science
studies: material entities as regards the behavior of the elemental
particles that constitute them, and the \emph{operations} it produces
consist in the acts of apprehending and judgment by the intellect
regarding that object. Note that without the \emph{habitus} of quantum
physics, such apprehensions and judgments would be impossible; thus, the
habitus is a true principle, or origin, of operation, albeit a secondary
one that depends on the intellect and ultimately on the substance (that
is, the human being) in which it inheres.

Similarly, if we look at an art---say, carpentry---the \emph{subject}
is, of course, the intellect; the \emph{object} is the works of
carpentry produced, which consists of pieces of wood to be shaped into
tables, chairs, and the like; and the operations consist in the
combination of knowhow and physical actions required to shape the wood
into the products desired. Table \ref{tab:subject-operation-object}
summarizes the examples we have just gone over, with some additions.

\begin{table}[ht]
\begin{center}
\caption{The subject, operations, and object of various faculties and operative \emph{habitus}.}
\label{tab:subject-operation-object}
{
\disablehyphenation
\begin{adjustbox}{max width=\textwidth}
\begin{tabular}{l|l|L{0.25\linewidth}|L{0.25\linewidth}}
\toprule
 & Subject & Operation & Object \\
\midrule
Sight & Substance (body) & Vision & \emph{Ens} as colored \\ 
Intellect & Substance (soul) & Apprehension, judgment & \emph{Ens} as intelligible \\
Will & Substance (soul) & Desire, love & \emph{Ens} as desirable \\
\midrule
Quantum physics & Intellect & Apprehension, judgment (of Q.P.) & Material \emph{ens} at subatomic scales \\
\midrule
Carpentry & Intellect & Actions to shape wood & Wood to be shaped \\
\midrule
Logic & Intellect & Apprehension, judgment, reasoning & Works of reason \\
\bottomrule
\end{tabular}
\end{adjustbox}
}
\end{center}
\end{table}

\subsection{Operations}\label{operations}

Of course, the effect of having faculties and \emph{habitus} is the
capacity for \emph{operation} (Greek: ἐνέργεια, \emph{enérgeia}), which
is simply the philosophical term for any actions that are produced by a
substance. As we have seen, operations always have their \emph{ultimate}
origin in the substance itself, but the faculty that produces them,
together with any \emph{habitus} that perfects it, can rightly be
considered the \emph{immediate} origins of that operation. In our
oft-cited example, the scientist who contemplates Schrödinger's equation
is rightly considered the author or origin of the act (operation) of
contemplating. However, he engages in this operation by the mediation of
his intellect, which has been perfected by the science of quantum
physics (an operative \emph{habitus}).

Figure \ref{fig:substance-operation} illustrates the hierarchical
relationship among substance, faculty, \emph{habitus}, and operation.
Operations can be produced entirely by the substance that acts, or else
the substance may receive the power to act from an extrinsic source. For
instance, an animal's capacity for locomotion is entirely intrinsic to
itself; however, if it is trained to perform certain actions by a human
master (say, for a dog to perform tricks or to sniff for contraband), it
is acting based on training that it has received extrinsically.
Similarly, a human being can do acts of knowledge and love under his own
power; however, he can only do acts of supernatural knowledge (faith)
and supernatural love (hope and charity) under the influence of grace,
which he receives from God. However, logic deals with the works of
reason, which are the result of operations that are entirely intrinsic
to man and do not require an extrinsic power. As shown in the diagram,
all operations that are proper---that is, intrinsic---to the substance
that produces them ultimately have their origin in the substance's
\emph{actus essendi} or \emph{act of being}. This is a principle that
you will learn about in your course on metaphysics: but for now, it is
sufficient to say that it is a metaphysical principle that is the source
of all proper actuality (i.e., all the actuality that does not come from
outside) in a substance. The \emph{essence} (\emph{essentia}) signifies
the substance inasmuch as it is a metaphysical principle, a capacity
that it fulfills in two ways: as the substrate for the accidents, and as
the potency receives the act of being. As such, the essence acts as a
limit to and a measure of the intensity of the act of being. The first
perfections to emanate from the substance are the active potencies
(\emph{potentiae activae}), which, as we saw, include the faculties of
sentient substances. Those faculties that entail a considerable freedom
of action---especially the rational faculties in man---are further
perfected by operative \emph{habitus}.

Putting all this together, all operations emanate from the substance's
\emph{actus essendi}. Those operations that are entirely intrinsic trace
all of their actuality to this act of being; others require the
assistance of an outside power. Regardless, operations are produced by
the mediation of substance's active powers. In sentient and rational
beings, some of these powers---called \emph{faculties}---are capable of
a further perfection that we call operative \emph{habitus}. Thus, a
substance's operations pass through, as it were, two and possibly three
filters: the essence itself (which already sets certain limits to what
kinds of operations are possible), the faculties that produce the
operations, and, where this is possible, the \emph{habitus} that perfect
the faculties.

\begin{figure}
\centering
\pandocbounded{\includegraphics[keepaspectratio]{images/substance-operation-graphic.pdf}}
\caption{An illustration how a substance produces operations.
\label{fig:substance-operation}}
\end{figure}

As a final note, every operation, evidently, produces an effect, which
we call its \emph{work} (\emph{operatum}) or \emph{product}. This work
can be extrinsic to the substance that has produced it---for instance,
the chairs and tables produced by the efforts of a carpenter---or else
it can be immanent---such as the internal effects of the actions of our
intellect and will. It is important to note (and I will reiterate this
point later on) that \emph{the works or products of an operation are
distinct from the operation itself}. When it comes to logic, we will see
that logic studies the \emph{works} of the intellect, not the
\emph{operations} of the intellect, which would pertain to philosophical
anthropology.

\section{The Operations and Works of the
Intellect}\label{the-operations-and-works-of-the-intellect}

Even though the object of logic is the \emph{works} of reason, it is
still good for us to have an overview of what the intellect does.
Fundamentally, the intellect (more specifically, what Aquinas calls the
``possible intellect'') is capable of accomplishing two kinds of
operation: it can discover \emph{what} something is, and it can discern
\emph{whether} something is the case. We say in more precise
philosophical terms, that the intellect can discover a thing's
\emph{quiddity} (from the Latin \emph{quid est}, ``What is it?''), and
it can make a judgment regarding that quiddity's \emph{being}---that is,
whether, in fact, that quiddity \emph{is} or \emph{is not}.

\subsection{Simple Apprehension}\label{simple-apprehension}

The species of intellectual operation that discerns \emph{what}
something is, we term the \emph{simple apprehension}. To illustrate how
this operation works, we can do a simple exercise. Suppose I attempt to
make the following statement:

\begin{quote}
Ice cream.
\end{quote}

\noindent This ``statement'' surely conveys some information to the
reader; ``ice cream'' is a idea that we all understand. However, these
words, by themselves, are insufficient to make an assessment regarding
an actual state of affairs. You have no idea whether I am saying that I
enjoy ice cream, or that chocolate ice cream is my favorite flavor, or
that the ice cream that I have served you is melting. With the words
``ice cream,'' I convey merely a \emph{notion}, an answer to the
question ``What is it?''

In our discussion above, we have already seen something that answers the
question ``What is it?'': namely, the substance, or more precisely, the
\emph{essence}---the ontological principle in a substance that gives it
its definition. Through a process that you will see in your other
courses, a faculty that Aquinas calls the agent intellect is able to
\emph{abstract} whatever is intelligible from the material entities that
we encounter. This intelligible notion or definition (Latin:
\emph{ratio}) is what actually answers the question ``What is it?'' in
our intellect, and thus we give it the technical name \emph{quiddity}
(Latin: \emph{quidditas}). It should be noted that the quiddity is not a
distinct ``thing'': it is not distinct from the object that is know.
Rather, it is is merely that object \emph{inasmuch as it is
intelligible} to us.

Ontologically speaking, only the substance can strictly be an
``essence''; that is, only a subsistent entity can answer the question
``What is it?'' and indeed, material substances are the objects that our
intellects most readily apprehend. However, we are \emph{also} able to
understand other entities, such as the accidents of substances, and
(through the use of negative judgments and analogy) even entities that
transcend the material world, or that can only be understood by analogy.
As a consequence, in \emph{logic} (but not in metaphysics), it is
necessary to expand our idea of ``quiddity'' to include notions that are
not substances, and even some notions that are not completely
unified---notions that are \emph{analogical} rather than
\emph{univocal}, as we will see in class.

Thus, the simple apprehension is what enables us to understand
\emph{quiddities}, and the product or work of that operation we call the
\emph{concept}\footnote{From the Latin \emph{conceptus}, that which has
  been conceived.}. As we will see in class, the concept can also be
called the \emph{universal} (technically, the \emph{universale in
praedicando}), because a single concept allows us to know the quiddity
of multiple individual things; and it can be called an
\emph{intention}\footnote{From the Latin roots \emph{in} (towards) and
  \emph{tendo} (to stretch out, to proceed, to strive).}, because it
leads us toward the thing that is known. Ontologically, a concept is a
``passion of the soul'' (Latin: \emph{passio animae})---that is, the
immanent effect on the soul of the act of apprehending. Since the
concept is one of the objects of our study, we will study it in much
greater detail in class.

\subsection{Judgement}\label{judgement}

Clearly, the simple apprehension by itself leaves us with very
incomplete knowledge: although it answers the question \emph{what}, but
itself, the apprehension tells us nothing about the actual state of
affairs. Going back to our example, I can use the concept ``ice cream''
to make a series of propositions like the following:

\begin{itemize}
\tightlist
\item
  This (what I have in my hand) is ice cream.
\item
  The ice cream is melting.
\item
  The ice cream has a chocolate flavor.
\item
  I ate ice cream yesterday.
\end{itemize}

\noindent Notice that in each statement, I am uniting two distinct
concepts\footnote{Observant readers will note that although each
  proposition has a single subject and a single predicate, some of the
  predicates are, in fact, \emph{compound}---that is, they unite two
  concepts into one. For example, if I were to draw out the last
  proposition into its elementary enunciations, they would be something
  like: ``I ate; the eating took place yesterday; ice cream was eaten.''
  Our language makes shortcuts that greatly simplify the expression of
  our judgments.}---or as Aristotle would say, I am \emph{composing}
them. That is, I am affirming that the concept on the right-hand side
(which we call the \emph{predicate}) truly \emph{is} a reality that
pertains to the concept on the left-hand side (which we call the
\emph{subject}). Subject and predicate, in this case, are united, either
explicitly (as in the first two statements) or implicitly (in the last
two) by the verb \emph{to be}, which we call the \emph{copula}. By
composing subject and predicate in this way, I am affirming the
\emph{being} of the predicate in the subject.

However, suppose that instead, I make the following statements:

\begin{itemize}
\tightlist
\item
  This (what I have in my hand) is not ice cream.
\item
  The ice cream is not melting.
\item
  The ice cream does not have a chocolate flavor.
\item
  I did not eat ice cream yesterday.
\end{itemize}

As before, I have taken two concepts, which can be termed \emph{subject}
and \emph{predicate}; however, this time, I have \emph{denied} the being
of the predicate. Logically speaking, instead of composing subject and
predicate, I have \emph{divided} or \emph{separated} them.

From the discussion above, I think it is clear why Aristotle and Aquinas
call this operation \emph{composition and division}, because it always
takes two concepts---a subject and a predicate---and either unites or
separates them. It can also be termed the ``judgment'' (and this was the
terminology preferred in the Modern period). Like the simple
apprehension, the judgment produces a work that is immanent to the soul,
which we call an \emph{enunciation}. Just as concepts are represented by
terms, enunciations are represented verbally as \emph{propositions},
like the statements I made above. It should be clear that it is only
after making a judgment and producing an enunciation that it makes sense
to talk about truth and falsehood. Those enunciations are true that
correspond to reality, and likewise, those that are in disagreement with
reality are false. That is, if I affirm something that actually exists,
then my affirmation is true; similarly, an enunciation denying something
that does not exist is also true. However, if I affirm something
nonexistent, or deny something that does exist, my enunciation is false.
Once again, it is the \emph{product} or \emph{work}, immanent in the
soul, that is the object of study of logic.

\subsection{Reasoning}\label{reasoning}

When we make a judgment regarding being, we can arrive at it in two
ways: either immediately, simply by seeing it, or by the mediation of
judgments that we have previously made. For example, any judgment that I
make regarding things that I can verify directly with my senses---say,
the fact that the pine tree outside my window is green---is immediate.
The intellect, by acts of apprehension, generates the concepts ``pine
tree'' and ``green,'' and by means of a judgment constructs a relation
of attribution between them, thus producing an enunciation: ``The pine
tree is green.'' In other words, producing that enunciation does not
require anything except the apprehension of those concepts and verifying
that they do, in fact, exist in reality. Suppose, however, that having
learned about the biology of pine trees, I make the following
enunciation: ``Pine trees are green because the chloroplasts in the pine
needles predominantly absorb all of the wavelengths found in visible
light \emph{except} green.''\footnote{Attentive readers will note that
  this enunciation, unlike the first one, is actually a \emph{compound}
  enunciation that includes a number of simpler ones joined together.}
There is no direct way to verify the existence of chloroplasts, still
less the fact that they tend to absorb red and blue light more than
green light. Reaching this conclusion requires highly specialized
instruments and a solid foundation in cell biology. In short, it is
impossible to make a judgment regarding the being of the chloroplasts
without extensive use of \emph{previously established enunciations}.

The discursive process that we use to make judgments that rely on
previously obtained knowledge is called \emph{reasoning} or
\emph{argumentation} (Latin: \emph{ratiocinium}). As we will see in the
course, argumentation entails constructing a series of \emph{arguments},
which consist of two previously known enunciations called
\emph{premises} that are linked by a concept in common called the
\emph{middle term}. By linking the two enunciations and dropping the
middle term, we produce a third enunciation called the
\emph{conclusion}. As with concepts and enunciations, arguments are
internal, intelligible passions of the soul that are distinct from their
external, verbal expressions, which we call \emph{syllogisms}. For
example,

\begin{quote}
All \emph{men} are mortal.\\
John is a \emph{man}.\\
Therefore, John is mortal.
\end{quote}

\noindent Of course reasoning made in real life will involve a great
many arguments, but every reasoning can be broken down to simple
arguments like this one, which, in its verbal expression, is called a
\emph{categorical syllogism}. Note that the intellect is not employing a
third operation, exactly; it is still using the second operation, the
judgment. What changes is merely the means that it uses to arrive at its
judgment on being: that is, previous knowledge, as opposed to direct
evidence.

\begin{table}[ht]
\begin{center}
\caption{A summary of the operations of the intellect, their work, and what they enable us to understand.}
\label{tab:operations-and-works-of-reason}
\begin{adjustbox}{max width=\textwidth}
\begin{tabular}{l@{\hskip -50pt}l@{\hskip 3pt}ll@{\hskip 3pt}l@{\hskip 3pt}l}
\toprule
Operation & & & Work & & What It Understands \\
\midrule
Simple apprehension & & & Concept & & \emph{What} something is (quiddity) \\
\multirow{2}{*}{Judgment} & \ldelim\{{2}{3mm} & immediate & Enunciation & \ldelim\}{2}{3mm} & \multirow{2}{*}{\emph{Whether} something is (being)} \\
 & & mediate & Argument & & \\
\bottomrule
\end{tabular}
\end{adjustbox}
\end{center}
\end{table}

Table \ref{tab:operations-and-works-of-reason} summarizes the role of
the two operations of the intellect and the three works they
produce.\footnote{According to some logic textbooks, there are
  \emph{three} operations of the intellect: apprehension, judgment, and
  reasoning. In my opinion, this is not entirely correct. It is better
  to say that there are two operations: apprehension and judgment.
  However, in order to arrive at some judgments, arguments must be
  employed.} These works---concepts, enunciations, and arguments---form
the basis of the three branches of logic: logic of the concept, logic of
the enunciation, and logic of the argument (or syllogism). At this
point, it bears reiterating: \emph{logic studies the works of reason,
not its operations}. What we will be discussing together are concepts,
enunciations, and arguments; it is not our task, in this course, to
consider apprehension and judgment, except insofar as these are the
\emph{cause} of the works. In fact, these works form the
\emph{object}---or, as we will learn to call it, the \emph{material
object} of logic. We will discuss the object of logic in detail in the
next chapter.

\chapter{The Definition of Logic}\label{the-definition-of-logic}

In this chapter, we will come to a detailed understanding of the
definition of logic, and thus we will finally be in a position to answer
the question \emph{What is logic?} definitively. I will go ahead and
state the definition here, but, of course, we will explain each part of
it in detail:

\begin{quote}
Logic is the instrumentally speculative and modally practical science
whose material object is the works of reason, whose formal object
\emph{quod} is the relations of reason constructed between second-order
intentions, and whose formal object \emph{quo} is the light of the first
principles of reason.
\end{quote}

We have already laid the necessary groundwork. In fact, you should
already be able to explain what it means for logic to be an
\emph{instrumentally speculative and modally practical}. By the end of
this chapter, you should have a good understanding of what logic
actually is.

\section{The Object}\label{the-object}

As we saw above, all faculties and operative \emph{habitus} have an
\emph{object}; that is, the entities that are referred to that faculty
or \emph{habitus}' operation. We gave a number of examples above, but we
need to look at the notion of object a bit more carefully. Let us take,
for example, the faculty of vision, which we said has as its object
colored bodies inasmuch as they are colored, or as Aquinas puts it,

\begin{quote}
Now, properly speaking, we call something the \emph{object} of a faculty
or \emph{habitus}, when all things are referred to that faculty or
\emph{habitus} by reason of that thing. For instance, \emph{man} and
\emph{stone} refer to sight insofar as they are \emph{colored}.
Therefore, \emph{colored things} are the proper object of
sight.\footnote{\emph{S.Th.}, I q.~1, a.~7, \emph{responsum} (my
  translation).}
\end{quote}

\noindent Notice how the object practically defines the faculty: vision
is the sense that allows us to see colored things. Let us unpack this a
little, because it will help us to understand how to describe the object
of an art or science like logic.

We said in the first place that the object of vision is colored material
bodies, or \emph{res colorata}. These bodies constitute, if you will,
the set of all things that can be the object of sight: it is, obviously,
impossible to see anything that is immaterial or invisible. When a
material entity becomes the object of our sight, it does so thanks to a
quality called \emph{color}: that is, an active power that allows the
object to emit a medium that our eyes can detect and that our brain can
interpret. Finally, there is the means that is used to generate color,
which is also the medium that transmits it to eyes: that is,
\emph{light}.\footnote{The ancients did not know that light is a kind of
  electromagnetic radiation, with a wavelength between that of infrared
  and ultraviolet radiation, that is, between 400 to 700 nanometers. As
  it turns out, unless you are viewing light emitted directly by a light
  source, the color of an object is produced when light from a light
  source (usually one that emits white light with a wide range of
  wavelengths) is absorbed by that object and is partly reflected into
  your eyes. The object tends to absorb some wavelengths better than
  others and reflect the rest. Thus, if the object is red, it is
  predominantly absorbing light with broadly blue and green hues, while
  reflecting back mostly red hues. Our eyes, in turn, fundamentally have
  three types of color detectors called \emph{cones}. Each type is most
  adept at detecting light of a particular hue: red, green or blue. By
  combining the signal detected by each kind of cone, the brain is able
  to construct the full range of colors that we are familiar with. None
  of this fundamentally changes the fact that there is a \emph{quality}
  of the object that is seen that is able to receive the incident light,
  that under the action of the light becomes an active power that then
  transmits the emitted light to our eyes. We call that quality
  \emph{color}.}

This brief analysis suggests that we may discern three aspects of the
object of sight: \emph{res colorata} constitutes what we will call the
\emph{material object} of the sense of sight; \emph{color}, which is the
aspect of colored bodies that our eyes are able to detect, we will call
the \emph{formal object quod};\footnote{In Latin, the \emph{obiectum
  formale quod}; literally, the ``formal object \emph{that}'' we see.}
and light, which is the means by which we see the object, we will call
the \emph{formal object quo}.\footnote{In Latin, the \emph{obiectum
  formale quo}; literally, the ``formal objet \emph{by which}'' we see.}

Let us unpack this further. As we saw in the last chapter, in
metaphysics \emph{matter} refers to a metaphysical principle that, by
itself, lacks definition or actualization, and requires an active
principle, a \emph{form}, in order to be something. We saw, for
instance, how substances function as a sort of secondary matter in which
the accidents are inherent (as when marble takes on the figure of David
or Apollo); and also how the human person is composed of a \emph{prime
matter}, the body, that is informed by a \emph{substantial form} that we
call the \emph{soul}.

When we use the terms \emph{material object} and \emph{formal object},
we are, of course, applying the metaphysical concepts of matter and form
by analogy. The material object is ``material'' in the sense that, by
itself, it lacks definition: it is not ``material'' because colored
bodies contain the ontological principle called prime matter (which, as
it happens, they do), but because it could not be presented to our sense
of sight without the particular aspect that \emph{makes} it apt for
vision. As we saw, that specific of \emph{res colorata} is precisely
\emph{color}. Because it gives the material object its definition, we
call color a \emph{formal} object, and we call it a formal object
\emph{quod} because color is \emph{that which} our sight actually sees.
We can go even further and specify the \emph{means} by which we see
color---namely, light---and since this functions as a kind of efficient
cause of our seeing, it is, as it were, even more ``formal'' or
determinative than the formal object \emph{quod}. For this reason we
call light the formal object \emph{quo}, because it is the means
\emph{by which} we see.

Let us look at another example that is, perhaps, closer to our goal: an
art like carpentry. As we saw previously, every art is a kind of
intellectual \emph{habitus} that permits the artist to perform certain
operations (not necessarily strictly intellectual in nature) that he
would be unable to do without possessing the art. In this case, the set
of objects that we are interested are the things that the carpenter will
actually work with; that is, things made of \emph{wood}. Thus, we can
say that the material object of carpentry is \emph{wood}. Of course, the
purpose of carpentry is to transform raw wood, as it is obtained in the
lumber yard, into finished products such as tables and chairs. Thus, we
could say that the formal object \emph{quod} of carpentry is wood
\emph{insofar as it is capable of being shaped}. Poor quality wood, or
wood that has already been made into a finished product, is most likely
\emph{not} something that the carpenter will use as his raw material
(unless perhaps he is fixing or restoring it). Finally, the formal
object \emph{quo} consists of the various techniques that the carpenter
uses in order to make his finished products.

When the \emph{habitus} in question is a science, the material object
will be what Aquinas calls the \emph{genus subiectum} and Aristotle τὸ
γένος τὸ ὑποκείμενον (\emph{tò génos tò hypokeímenon}): the
\emph{subject matter}, or the set of things that the science studies.
Thus, for example, biology studies \emph{living things}, and these
constitute its material object. The formal object \emph{quod} of a
science will be the particular aspect of that material object that is of
interest to the science. For instance, biology does not study living
things in every respect---it does not, for example, look in to the
metaphysical principles of living things---but rather only it studies
them \emph{insofar as they are alive}. Thus we could say that the formal
object \emph{quod} of biology is the \emph{vegetative and animal life}
or \emph{bodily life} of the living things it studies. Finally, the
formal object \emph{quo} of a science will be the method used: in the
case of biology, the scientific method, which entails observation,
forming hypotheses, making deductions, and conducting experiments and
other date-gather activities. Table \ref{tab:examples-of-objects}
summarizes these examples.

\begin{table}[ht]
\begin{center}
\caption{Examples of the objects of various faculties and \emph{habitus}. Note that I have listed some faculties first, followed by arts, and finally some sciences.}
\label{tab:examples-of-objects}
{
\disablehyphenation
\begin{adjustbox}{max width=\textwidth}
\begin{tabular}{l|l|l|l}
\toprule
 & Material Object & Formal Object \emph{Quod} & Formal Object \emph{Quo} \\
\midrule
Sight & Colored \emph{Ens} & Color & Light \\
Intellect & \emph{Ens} & Intelligibility & The first principles of reason \\
Will & \emph{Ens} & Desirability & The first principles of morality \\
\midrule
Carpentry & Workable wood & Suitability for being shaped & Woodworking techniques \\
Painting & Things that can be depicted & Ability to be depicted & Painting techniques \\
\midrule
Biology & Living things & Material life & Scientific method \\
Quantum physics & Material \emph{ens} & Subatomic scale & Scientific method \\
Metaphysics & \emph{Ens} & \emph{Ens} & Resolutive reasoning \\
\bottomrule
\end{tabular}
\end{adjustbox}
}
\end{center}
\end{table}

You will often see a science described using a formula such as the
following: ``Science \emph{S} studies entity \emph{E} inasmuch as it has
property \emph{X}.'' For example, ``Metaphysics is the study of entities
(\emph{ens}) inasmuch as they \emph{are} (\emph{qua ens}),'' or,
``Biology is the study of living things inasmuch as they have bodily
life.'' This is simply a short way to define a science in terms of its
material object (entity \emph{E}) and formal object \emph{quod}
(``inasmuch as it has property \emph{X}'').

One final observation: there is \emph{one} object of a given faculty or
\emph{habitus}. The language that we have used here might give the
impression that there are three objects, but in reality these are three
\emph{aspects} of a single object: the things that we have identified as
referred to the faculty or \emph{habitus}. Thus, to reiterate, these
things are \emph{materially speaking} some sort of entity (real or in
our reason alone, as we will see) and we call them the \emph{material
object}. Formally speaking there is some specific \emph{aspect} that
refers these entities to the faculty or \emph{habitus}, which we call
the \emph{formal object quod}. Finally, there is a means that permits
the referral to take place, which we call the \emph{formal object quo}.

\section{The Material Object of
Logic}\label{the-material-object-of-logic}

We must now, of course, turn our attention to logic itself. What is the
material object of logic? The answer, fortunately is relatively
straightforward. Logic studies the \emph{works of reason}, which I have
outlined in the previous chapter: concepts, enunciations, and arguments.
These, as we saw, are \emph{passiones animae}, that is, qualities of the
soul that remain immanent there, due to the operations of the intellect
(which, as we saw are the \emph{apprehension} and the \emph{judgment},
the latter operating either immediately or mediately).

I will take a moment to emphasize once again: \emph{logic studies the
works of reason, not the operations of reason}. We are interested in
forming precise and well defined \emph{concepts}, \emph{enunciations}
that are properly formed and expressed, and \emph{arguments} that are
both correctly constructed and true in their conclusions. In logic, the
manner in which the intellect functions does not enter the object as
such: we are only interested in the intellect's operations as the
\emph{causes} of that object.

We must, however, now tackle the much more difficult issue of the
\emph{formal object quod} of logic: what is it about the works of reason
that logic concerns itself with? What \emph{aspect} of these works does
it focus on?

\section{\texorpdfstring{The Formal Object \emph{Quod} of
Logic}{The Formal Object Quod of Logic}}\label{the-formal-object-quod-of-logic}

As the definition I gave above states, the formal object \emph{quod} of
logic regards \emph{relations of reason} that we construct between
second-order intentions. Before we get any further, we will need to
discuss \emph{relations}, \emph{relations of reason}, and
\emph{intentions}.

\subsection{Relations}\label{relations}

Consider for a moment a father and a son. Metaphysically speaking, both
are, of course, substances. However, once the father has begotten a son,
he acquires a new perfection that he did not have before:
\emph{fatherhood}. Likewise, the son simultaneously receives a
perfection upon his conception called \emph{filiation} or
\emph{sonship}. Those perfections are not, themselves, substances, but
they are real attributes; hence, they must be what we have learned to
call \emph{accidents} in the last chapter. In fact, Aristotle considers
perfections such as these to be their own category, or supreme genus, of
accident that he calls the πρός τι or ``toward which.''

In short a πρός τι, or, as we have come to know it through Medieval
Latin, a \emph{relation} is, in the first instance, a reference that one
substance has to another.\footnote{For anyone who has studied Latin, it
  can be helpful to recall that the word \emph{relation} comes from the
  verb \emph{re-ferro}, which takes the past participle \emph{re-latum}.}
You will learn more about relations in your logic course, since they
constitute one of the categories that we will study in detail. As we are
about to find out, the concept of \emph{relation} is extended beyond
that of the real perfection that is inherent in a substance. We
sometimes call this kind of relation a \emph{real relation} so as to
distinguish it from others that we are about to look at. Table
\ref{tab:examples-of-real-relations} gives some examples of real
relations.

\begin{table}[ht]
\begin{center}
\caption{Some examples of real relations. Note that relations are always conceived \emph{bilaterally}: that is, the relation from the subject to the object will normally be coupled to a relation from the object to the subject. However, in some cases, the reciprocal relation is not a real one. The relation “having known” is real; the reciprocal relation, “being known,” is not.}
\label{tab:examples-of-real-relations}
{
\disablehyphenation
\begin{adjustbox}{max width=\textwidth}
\begin{tabular}{l|l|l|l}
\toprule
Relation & Subject & Object & Foundation \\
\midrule
Fatherhood & the father & the son & generation \\
Sonship & the son & the father & generation \\
Double & a two-meter-long substance & a one-meter-long substance & quantity \\
Half & a one-meter-long substance & a two-meter-long substance & quantity \\
Having known & the knower & the think known & intentional similarity \\
\bottomrule
\end{tabular}
\end{adjustbox}
}
\end{center}
\end{table}

As I said, we will see this in detail in the course. For now, a few
observations about relations are all that we need. First, some
terminology: the substance in which the relation is inherent we call the
\emph{subject} of the relation (as we would for any accident). The
substance referred to we call the \emph{object}. Moreover, the relation
will always be based on some kind of \emph{foundation} (quantity, action
and passion, or some kind of knowledge). Second, real relations are
normally \emph{bilateral} or \emph{reciprocal}. That is, the relation
inherent in a subject will ordinarily be paired with \emph{another}
relation that is inherent in the substance that is the object of the
first relation. The example we gave above of fatherhood and sonship
illustrates this well: a father cannot be a father unless he has a child
(in our example, a son), and that son must necessarily have a father.
For this reason, relations are always \emph{conceived} as bilateral ``by
default'' by our intellect---even when it turns out that the reciprocal
relation does not really exist. (We will see some examples below.)
Finally, in modern parlance, we tend to refer to the \emph{pair} of
relations as a \emph{relationship}. For Aristotle and Aquinas, however,
what we call a \emph{relationship} actually consists in two distinct
\emph{relations}, each one a perfection inherent in a distinct subject.
Thus, for example, in modern parlance we would say that a husband and
wife have an indissoluble \emph{relationship} that we call the
\emph{marriage bond}; however, in metaphysical terms, that bond consists
in \emph{two} relations, one inherent in the husband, and the other
inherent in the wife. A similar things can be said, say, of the
father--son relationship and the teacher--student relationship.

\subsection{Relations of Reason}\label{relations-of-reason}

I have brought up real relations at this point not because they are
directly pertinent to our topic, but because they are necessary for
understanding an essential element of the definition we are considering:
\emph{relations of reason}.

A \emph{relation of reason} (\emph{relatio rationis}) can be though of
as a construct that our reason treats \emph{as if} it were a real
relation, when in reality it is nothing at all. A few examples will
illustrate the idea. Consider the following proposition:

\begin{quote}
John is himself.
\end{quote}

\noindent What we have done here is subtle. The subject of the
proposition is \emph{John}, but the predicate is also \emph{John}: the
term \emph{itself} is merely a placeholder to avoid being repetitive.
Thus, \emph{John} is, in a sense ``divided'' by our reason; that is, we
consider two ``instances''~of him and attribute to each one a kind of
relation. In effect we are saying that John (in the subject) is
\emph{the same} as John (in the predicate), and vice versa. However,
this relation, which we can call \emph{sameness} or
\emph{identity}\footnote{From the Latin \emph{idem}, meaning ``the
  same''. Readers preparing to write papers should be familiar with the
  Latin term \emph{ibidem} (from \emph{ibi}, ``there'', and
  \emph{idem}), which means ``in the same place.''}, does not actually
exist. John (in the subject) is not really distinct or different from
John (in the predicate). The relation of \emph{identity}, therefore,
does not really exist, but rather is merely a \emph{construct} to help
us understand that substances are united and permanent---that is, a
substance (or for that matter, any entity) cannot simply become
something else without some kind of real transformation.
\emph{Identity}, therefore is a \emph{relation of reason}, not a real
relation.

Or else, consider real relations that are based on knowledge. Suppose,
for example, that you have read a book---say an introductory textbook on
logic. Now that you have read it, the relation of \emph{having read} the
book is now inherent in you. However, the book, as such, is completely
unchanged by you having read it. As I mentioned earlier, we are unable
to conceive of a relation in any other but as bilateral or reciprocal,
and so our reason cannot help but consider the relation (ostensibly
inherent in the book) of \emph{having been read} alongside the real
relation. In other words, when you read a book, you acquire a new
perfection, \emph{having read} the book; the book acquires nothing at
all, but we cannot help but conceive a \emph{having been read} in tandem
with \emph{having read}. However, we use our judgment to deny that
\emph{having been read} actually exists: it is, in short, a
\emph{relation of reason}.

In order to understand these relations of reason, it helpful to mention
that they belong to a broader group of pure objects, or constructs, that
we call \emph{entia rationis}, entities of reason. An \emph{ens
rationis} refers to any construct that we treat \emph{as if} it were an
entity, when in reality it is not. It is, if you will, a placeholder
concept which, upon examination, does not refer to any actual
\emph{thing} (\emph{res}) in reality.

As usual, it is easiest to understand the idea using examples. Consider
\emph{hobbits}. If anyone is not familiar with the idea, hobbits are a
fictional people imagined by J.R.R.~Tolkien in his fantasy works,
\emph{The Hobbit} and \emph{The Lord of the Rings}. In reality there
\emph{are} no hobbits. They are not, in the strict sense \emph{essences}
or \emph{substances}; however through our imagination, we can describe
what such a people would be like if they \emph{did} exist. In other
words, hobbits are \emph{figments}, with no actual ontological
foundation. In fact, this analysis is valid for any fictional character
or object: Greek gods, superheroes, Romeo and Juliet, light sabres, and
so on.

On the other hand consider a construct such as \emph{darkness}. Darkness
is not really an entity; it is best characterized as the \emph{lack} of
light. It is, however, clearly based on real situations: darkness does
not exist as such, but dark \emph{rooms}, dark \emph{nights}, and dark
\emph{clothing} do. A similar thing can be said of other privations:
\emph{coldness}, \emph{evil}, \emph{blindness}, and the like. In each of
these cases, we are, as it were, hypostatizing the \emph{lack} of a real
thing: \emph{coldness} is the lack of heat; \emph{evil} is the privation
of good;\footnote{It might not be immediately obvious that evil is a
  privation, but if you think about it, there are only evil
  \emph{persons}, evil \emph{circumstances}, and evil
  \emph{actions}---that is, real entities that have somehow been
  corrupted. However, evil, in and of itself, does not exist.}
\emph{blindness} is the privation of sight in something that ought to be
able to see. There are also the relations of reason that I mentioned
above, such as \emph{identity} and \emph{having been read}.

A common thread among all the \emph{entia rationis} in this second group
is that although each one is a pure object that, as such, does not
exist, it nevertheless has a very precise \emph{foundation} in reality,
one that is entirely lacking in the first group. These examples suggest
a twofold division of \emph{ens rationis}: one according to the how the
\emph{ens rationis} is attributed to a subject, and the the other
according to the way the \emph{ens rationis} is constituted as a pure
object of reason.

Let us unpack this idea. When I use \emph{hobbit} as a predicate---as in
the proposition ``Frodo Baggins is a hobbit,'' the subject to which I am
attributing the concept \emph{hobbit} is completely nonexistent. There
is, and never was, a hobbit named Frodo Baggins. On the other hand, when
I say, ``The room is dark,'' the subject, \emph{room} is quite
real---even though the attribute, as such, is not. We can thus divide
\emph{entia rationis} into those that are attributed to completely
nonexistent subjects, and those whose subjects are real. We call former
class \emph{ens rationis sine fundamento in re} (the entity of reason
without any foundation in reality), and the latter \emph{ens rationis
cum fundamento in re} (the entity of reason having some foundation in
reality).

If we look only at \emph{entia rationis cum fundamento in re}, our
examples suggest that there are two ways to form them: either positively
(that is, based on an affirmation), or negatively (based on a negation).
It is not hard to see, based on the examples I mentioned, that the
negation of some real perfection, as such, is not truly a ``something,''
but rather the \emph{lack} of something. As regards those \emph{entia
rationis} formed by an affirmation, Aquinas argues that when we affirm
something in an absolute sense, that affirmation must refer to something
in reality: as when we say, ``The sky is blue,'' or ``Gold has a density
of 19.3 g/cm\textsuperscript{3}.''\footnote{There is a difficulty in
  interpreting the passage that talks about this distinction, which is
  in \emph{Quaestiones Disputatae de Veritate}: ``Now, what exists in
  reason alone {[}\emph{id quod est rationis tantum}, that is, \emph{ens
  rationis}{]} can only be of two kinds: that is, either a negation or a
  kind of relation. For every absolute affirmation signifies something
  found in the nature of things {[}i.e., reality{]}. Thus to \emph{ens},
  `oneness' adds something that is in reason alone, that is, a negation:
  an \emph{ens} is said to be `one' insofar as it is undivided. However
  `truth' and `goodness' are attributed positively; therefore, they
  cannot add anything but a relation that is in reason alone''
  (\textsc{Thomas Aquinas}, \emph{Quaestiones disputate de Veritate},
  Leon. vol.~22, Sanctae Sabinae, Rome 1975--1970--1972--1973--1976
  {[}henceforth cited as \emph{QdV}{]} 21, 1, \emph{corpus}; my
  translation).

  In this passage, Aquinas is talking about the \emph{trascendentia} or
  the ``properties'' of \emph{ens}, which later philosophers call the
  \emph{transcendentals}. We will study these in the logic of the
  concept, but the relevant connection to \emph{ens rationis} is that
  the notions conceptually distinguishing, say, oneness, truth, and
  goodness from \emph{ens} are \emph{entia rationis}---that is, they are
  not real things \emph{in rerum natura} (as Aquinas would say).
  However, those distinguishing notions do have a real foundation: they
  are attributed to real things, namely, all \emph{entia}. Aquinas'
  argument poses a difficulty: does he intend divide all \emph{entia
  rationis} without exception into negations and relations? If so, what
  are we to do with \emph{entia rationis sine fundamento in re}---those
  \emph{entia rationis} with no foundation in reality, such as hobbits
  and Greek deities?

  There are fundamentally two possible solutions: either \emph{entia
  rationis sine fundamento in re} can in some way conceived as relations
  of reason, or else he is not considering them in this division. I tend
  to favor the latter interpretation: in short, I think that our reason
  treats \emph{entia rationis sine fundamento in re} in a hypothetical
  way, \emph{as if} they were posited absolutely---and then, with the
  judgment, denies that they actually exist. A reader of \emph{The
  Hobbit} asks himself, ``What would Middle Earth have been like
  \emph{if} it existed and a hobbit named Bilbo Baggins lived there?''
  Thus, I think such complete figments are excluded from the division
  that Aquinas is considering here.} The only way to obtain an \emph{ens
rationis} through positive affirmation is to affirm it
\emph{relatively}---that is, as a relation of reason, as in the example
of the relations of \emph{identity} and \emph{having been read} that I
mentioned above. Table \ref{tab:ens-rationis} summarizes the divisions
of \emph{ens rationis}.

\begin{table}[ht]
\begin{center}
\caption{A summary of the divisions of \emph{ens rationis}.}
\label{tab:ens-rationis}
\begin{adjustbox}{max width=\textwidth}
\begin{tabular}{llll}
\toprule
& Division by attribution & Division by object & Example \\
\midrule
\ldelim\{{3}{*}[Ens rationis\hspace{12pt}] & sine fundamento in re & & hobbits \\
 & \ldelim\{{2}{*}[cum fundamento in re\hspace{12pt}] & negation & blindness \\
 &  & relation & having been read \\
\bottomrule
\end{tabular}
\end{adjustbox}
\end{center}
\end{table}

Before moving on, I will make one final note about the term \emph{ens
rationis}. I focused on the meaning used most often in philosophy---that
is, a nonexistent construct that our reason treats \emph{as if} it were
a thing in reality---however the term itself is ambiguous. The term
``entity of reason,'' according to its verbal construction, simply means
an entity that depends on reason in some way. Thus, in theory, it could
have three possible meanings:

\begin{enumerate}
\def\labelenumi{\arabic{enumi}.}
\item
  It could refer to the \emph{real} entities that have have reason as
  their \emph{efficient cause}: that is, it could be a synonym for the
  works of reason, namely, concepts, enunciations, and arguments. In
  that case, we say that their dependence on reason is \emph{effective}.
\item
  It could refer to the \emph{real} entities that have reason as their
  subject or substrate: that is, all of the intellectual \emph{habitus},
  such as sciences, arts, prudence, and the theological virtue of faith.
  In that case, we say that their dependence on reason is
  \emph{subjective}.
\item
  It could refer to the \emph{nonexistent} entities---in reality, pure
  objects or constructs---that our reason apprehends \emph{as if} they
  were entities. In this case we say that their dependence on reason is
  \emph{objective}.
\end{enumerate}

\noindent It goes without saying that all of the \emph{entia rationis}
that we have talked about in this section are of the third kind, those
that depend on reason \emph{objectively}---that is as pure objects or
constructs with no ontological consistency whatsoever.

With all of this background out of the way, the question remains: what
are the relations of reason that interest us in logic? In essence, they
are the relations that our reason constructs from one work of reason to
another: from a concept to a concept, and from an enunciation to an
enunciation. Here are a few examples that you will study during the
course:

\begin{itemize}
\tightlist
\item
  Attribution or predication: the relation by which we affirm (or deny)
  the being of the predicate in the subject of an enunciation.
\item
  Being the subject to which a predicate is attributed.
\item
  Predicability: the five ways in which a predicate can be attributed to
  a subject.
\item
  The relation of a species to the genus that contains it, and vice
  versa.
\item
  The relations of the major and minor premise to the conclusion, and
  vice versa, which constitute the logical nexus.
\end{itemize}

\noindent Since all of these relations have works of reason (which are
accidents, \emph{passiones animae}) as their subject and object, they
are clearly not \emph{real} relations, but relations of reason. These
relations will be the focus of our study throughout the course, and
identifying or constructing them \emph{correctly} will be our goal.

\subsection{Intentions}\label{intentions}

When you hear the word \emph{intention}, you are probably used to the
most common modern usage of the term, which references the movement of
the will: one \emph{intends} to go shopping if one has \emph{plans} to
go to the store later on, we apologize for hurting someone
\emph{unintentionally} (i.e., without willing it), we make prayer
\emph{intentions} (i.e., express our intent to pray for someone), and so
on. As I mentioned, in all of these cases there is a movement of the
will toward some object. \emph{Intention}, as you can probably tell,
comes from the verb \emph{intendō}, from the roots \emph{in} (in this
case meaning, ``toward'') and \emph{tendō}, meaning ``to stretch,
proceed, or strive for.'' Among several meanings, \emph{intendō} refers
to a movement toward something: to turn one's attention to something, or
(as in modern language) to plan to do something.

In Aquinas, \emph{intendō}, and the related term \emph{intentio}, can
refer to movements of the will, but it can \emph{also} refer to the
operations of the intellect. In particular, it can refer to the
\emph{concept} insofar as it is directed to an \emph{object}. For
example, the concept ``man'' represents all human beings, and in that
sense is directed toward (\emph{intendit}) them.\footnote{Attentive
  readers might realize that there is, in fact, a \emph{relation of
  reason} that refers the concept to its object. We call it
  ``intentional similarity.'' The \emph{real} relation inherent in the
  person who knows something through this concept we call
  \emph{intentionality}.} For this reason, Aquinas uses the term
\emph{intentio} sometimes as a synonym for the concept, not just as
referring to the movement of the will.\footnote{In fact, \emph{intentio}
  in Aquinas can refer to the direction of anything toward its end: it
  is, in this sense, the ``intention'' (in modern language,
  ``tendency'') of the well aimed arrow to reach its target. See, for
  instance, I-II q.~12, a.~2, resp.: ``intentio, sicut ipsum nomen
  sonat, significat in aliquid tendere.''}

Now, it is possible to form a concept (that is, an intention) of
\emph{real entities} (which is what we normally do), but it is also
possible to form a concept of \emph{the works of our reason}. For
example, when I make the statement, ``All human beings are mortal,'' the
term \emph{human being} there is referring to real flesh-and-blood human
beings. On the other hand, I can also look reflexively on the works of
reason that permitted me to make the statement: in this case, the
concepts \emph{human being} and \emph{mortal}, and the enunciation that
composed the two concepts together.

We call concepts that refer \emph{immediately} to real entities
\emph{first-order intentions}. On the other hand, concepts that refer to
some work of our reason, we call \emph{second-order intentions}. Note
that even second-order intentions refer to reality in a \emph{mediate}
or indirect way: for instance, when I take the term ``human being'' to
mean the \emph{concept} formed in my reason, that concept refers in the
first place to the \emph{concept}---but of course that concept, in turn,
refers to real human beings. I will let Aquinas speak here:

\begin{quote}
Something in reality can correspond to a concept in two ways. In the
first way, it does so immediately, for instance when the intellect
conceives the form of some thing that is found outside the soul, such as
a man or a stone. In the other way, it does so mediately, for instance
when something follows the action of the intellect, and the intellect
reflecting upon itself considers it. In that case, the reality
corresponds to that consideration of the intellect in a mediate way,
that is, by means of the understanding of reality. For example, the
intellect understands the nature of animals in a man, in a horse, and in
many other species. From this, it follows that the intellect understands
that nature as a genus. Nothing in external reality immediately
corresponds to the concept by which the intellect understands the genus.
However, some reality does correspond to the understanding from which
this intention flows.\footnote{\textsc{Thomas Aquinas},
  \emph{Quaestiones disputate de potentia}, in: \emph{Quaestiones
  disputatae}, vol.~2, Marietti, Turin 1965 {[}henceforth citied as
  \emph{QdP}{]} q.~1, a.~1 ad 10 (my translation).}
\end{quote}

Perhaps an analogy that can help is to compare the intellect with an
instrument like a telescope. Just as the telescope is the instrument by
which we are able to see faraway objects like stars, the intellect is
the instrument by which we are able to know things. In order for us to
see the stars, the telescope must focus the light using lenses and
similar equipment. In a way, our concepts play a role analogous to that
of the lenses: they are, if you will, the ``window'' through which we
see the realities that are outside our soul. So long as the lens is
focusing light coming from \emph{stars}, it is similar to a
\emph{first-order intention}: that is, it is helping us to see, in an
immediate way, the \emph{stars}---an external reality. At this point, we
will need to bend the analogy a little. Suppose that we could set up the
telescope in such a way that we could mount a lens capable of focusing
on \emph{another lens}---perhaps in order to observe the first lens in
action. That is how a second-order intention works: it reflexively looks
back on the intellect, and points to \emph{another work of reason}
rather than an reality outside the soul.

As usual, the distinction between first-order and second-order intention
will be much easier to understand with examples. Consider the following
two propositions:

\begin{itemize}
\tightlist
\item
  All men are mortal
\item
  ``Man'' is a species of the genus ``animal.''
\end{itemize}

\indent In the first proposition, the concept \emph{man} is a
first-order intention, because it refers to flesh-and-blood men. In the
second proposition, \emph{man} is a second-order intention, because it
signifies, not flesh-and-blood men, but the \emph{concept itself}.

In summary: a first-order intention is concept that refers directly to a
reality outside the soul. A second-order intention is an intention that
refers to another intention (or, at any rate, to some work of reason).

\subsection{Relations of Reason between Second-order
Intentions}\label{relations-of-reason-between-second-order-intentions}

We are finally in a position to understand the formal object \emph{quod}
of logic:

\begin{quote}
the relations of reason constructed between second-order intentions
\end{quote}

\noindent In short, \emph{logic in itself does not study realities that
are outside the soul}. Instead, it \emph{studies the works of reason as
such}, and as a consequence, we will be employing second-order
intentions. Logic does not concern itself with the density of
gold---that is for chemistry. However, it does concern itself with the
fact that in the proposition, ``Gold has a density of 19.3
g/cm\textsuperscript{3},'' the concept \emph{gold} is a species, it
functions as the subject of the enunciation, and a predicate (density of
a certain magnitude) is attributed to it as a property. As the example
demonstrates, logic deals specifically with the relations of reason that
the intellect constructs from one second-order intention to another
(e.g., the relation of reason that makes \emph{gold} the subject with
respect to \emph{density of 19.3 g/cm\textsuperscript{3}}, and the
reciprocal relation that makes \emph{density of 19.3
g/cm\textsuperscript{3}} the predicate of \emph{gold}).

What logical relations will we study? They are fundamentally of three
kinds or genera, which correspond to the three types of works of reason:

\begin{enumerate}
\def\labelenumi{\arabic{enumi}.}
\item
  \emph{Predicability}, which describes \emph{concepts}, and concerns
  the ways in which a concept can be attributed to its subject. As we
  will see, concepts can be either \emph{univocal} (with one well
  defined meaning) or \emph{analogical} (having only a partial unity).
  Moreover, the univocal concept is predicable on one of five ways: as
  \emph{species}, \emph{genus}, \emph{difference}, \emph{property}, or
  \emph{accident}.
\item
  \emph{Predication}, which concerns actual attribution of a predicate
  to a subject, thus constituting an \emph{enunciation}. We will learn
  to divide enunciations according to their mode of attribution by
  \emph{quantity}, \emph{quality}, \emph{unity}, and \emph{matter}.
\item
  \emph{Illation},\footnote{Note that ``illation'' comes from the Latin
    verb \emph{infero} (past participle: \emph{illatum}), from which the
    word \emph{inference} is derived.} which concerns the logical
  consequence of a conclusion from its premises, thus constituting an
  \emph{argument}. Illation can, in turn, be \emph{deductive}
  (proceeding from general premises to more specific conclusions) or
  \emph{inductive} (proceeding from specific premises to more general
  conclusions). We shall spend the majority of our study of the logic of
  argumentation on deductive reasoning, since this forms the backbone of
  all scientific argumentation.\footnote{As we will learn, we cannot
    arrive at certain, universal knowledge by relying solely on singular
    data. Pure induction, therefore, only produces \emph{probable}
    conclusions. However, contact with even one material entity gives
    some knowledge about that thing's \emph{nature} or \emph{quiddity}.
    That knowledge is sufficient to make universal enunciations about
    the genus or species to which that individual belongs. Thus, we are
    capable of true \emph{scientia} even though we gather our data from
    individuals.}
\end{enumerate}

\section{\texorpdfstring{The Formal Object \emph{Quo} of
Logic}{The Formal Object Quo of Logic}}\label{the-formal-object-quo-of-logic}

We have discussed the material object and formal object \emph{quod} of
logic: that is, \emph{what} logic studies---namely, the works of
reason---and \emph{what specific aspect} of the material object it is
concerned with---namely, the second-intention relations of reason that
allow us to know the truth with certainty. We must now discuss
\emph{how} logic is able give us the ability to construct these
relations correctly and thus bring us unfailingly to the truth.

As you will learn in your epistemology class, the first concept that the
intellect makes is that of \emph{ens}: that is as soon as the intellect
is active, it realizes that there is \emph{something} outside itself,
and that this something \emph{is}. We do not start out understanding
\emph{ens} in a formal or explicit sense: systematic study of
metaphysics is required for that.\footnote{If our intellection began
  with explicit and formal knowledge of \emph{ens}, we would expect
  children just learning to speak to turn to their mothers and say
  ``Entity!'' instead of ``Mommy!'' Since that is obviously not the
  case, we know that our first knowledge of \emph{ens} is implicit.}
Instead, this knowledge is implicit, acquired by the very fact that our
intellect is working: \emph{actu exercito} (through the act that is
exercised), as Medieval philosophers would put it.\footnote{We say that
  knowledge gained formally and explicitly is obtained \emph{actu
  signato}, through an act (of apprehension or judgment) that is
  explicitly designated for this purpose.}

As soon as we acquire our first concept, \emph{ens}, our agent intellect
causes us to form our first enunciation at the very same moment: namely
that \emph{ens} is not \emph{non-ens}, in other words, whatever it is
that is presenting itself to the intellect must, on some level,
\emph{be}, and cannot possibly \emph{not be}, since it is presenting
itself right now. If, for example, I were to get up in the middle of the
night and stub my toe on some hard object, the first thing I would know
is that my toe has struck \emph{something}, some entity, even if I do
not immediately know what it is. It is also clearly not \emph{nothing}:
that object, which is presenting itself painfully to me is clearly
something that \emph{is}, and it cannot possibly \emph{not} be. This
judgment that we make, also \emph{actu exercito}, whenever we are
presented with an entity is called the \emph{principle of
non-contradiction}, the very first principle of reason.

As it turns out, the very exercise of our intellect reveals a whole
store of first principles: the principle of causality (namely, that
every entity presenting itself to us is on some level an effect that
must have a cause), the principle that the whole is greater than the
part, the principle of identity (that an entity is itself, and not some
other entity), and so on. Aquinas argues that this set of principles
that our reason, even before our formal study of logic, infallibly sets
forth constitutes an \emph{habitus} that he calls the \emph{habitus
primorum principiorum.}\footnote{Rather confusingly, Aquinas also gives
  this \emph{habitus} the name of \emph{intellectus}. Thus, it whenever
  the term \emph{intellectus} appears in his works, it is very important
  to understand which meaning he is giving the term at a given moment.}
Since this \emph{habitus} derives its efficacy from the very exercise of
our reason, it is present in all human beings, infallibly, and without
exception.

It is these principles that make logic possible. In fact, if you recall,
I said that it was defensible to call reason itself a kind of ``natural
logic.'' This is true because of the first principles of reason.

\section{Putting All the Pieces
Together}\label{putting-all-the-pieces-together}

Therefore, to summarize what we have learned so far:

\begin{itemize}
\item
  Logic is an instrumentally speculative and modally practical science.
  We saw in the first chapter how logic fulfills the characteristics of
  both a \emph{science} and \emph{art}, and, indeed can be characterized
  as both a speculative and a practical science. It can be all three at
  once because, unique among all intellectual \emph{habitus}, its very
  object of study, the works of reason, is also what it produces, and it
  is also the means by which it reaches the truth. Because logic is the
  means par excellence by which all other sciences attain true knowledge
  of the world outside the soul, but does not seek that kind of
  knowledge in itself we say that it is \emph{instrumentally
  speculative}. Because logic helps other sciences obtain that knowledge
  by the correct construction of relations between the works of reason,
  we say that logic is \emph{modally practical}.
\item
  The material object of logic is the works of reason. Logic concerns
  itself as such with concepts, enunciations, and arguments: the fruit
  of the operations of our intellect. It bears reiterating: \emph{logic
  studies the }works* of reason, not the operations of reason*.
\item
  The formal object \emph{quod} of logic---that specific aspect of the
  material object that logic is concerned with---is the relations of
  reason that can be constructed between second-order intentions. Logic
  looks reflexively on the works of reason, and tells us how they are to
  be ordered if we are to acquire true and certain knowledge. It is all
  the other sciences---metaphysics, mathematics, physics, chemistry,
  biology, economics, astronomy, and so on---that deal with \emph{first}
  intentions.
\item
  Finally, what permits us to use our reason in a logical fashion (what
  we called natural logic), and then to study the works of reason (what
  we called scientific logic), is the first principles of natural
  reason. Thus, we say that the formal object \emph{quo} of logic is the
  light of the first principles.
\end{itemize}

In short,

\begin{quote}
Logic is the instrumentally speculative and modally practical science
whose material object is the works of reason, whose formal object
\emph{quod} is the second-intention relations of reason constructed
between those works, and whose formal object \emph{quo} is the light of
the first principles of reason.
\end{quote}

\section{The Divisions of Logic}\label{the-divisions-of-logic}

Now that we have established what logic actually is, I can say something
about how the course is organized. The works of reason---concepts,
enunciations, and arguments---constitute the material object of logic,
as we saw, and for this reason we divide classical logic into three main
chapters: the logic of the concept, the logic of the enunciation, and
the logic of the argument or syllogism.

Getting to the truth is always the ultimate goal of logic. However, as
we will see in the course, a large part of our work in logic involves
what we will call understanding the \emph{formal conditions} for
obtaining the truth: that is, if we presuppose that we are already
starting with sound principles and true knowledge, we can understand how
to proceed to \emph{new} knowledge starting from the old. We call this
aspect of logic \emph{formal logic} or \emph{minor logic}: the part of
logic that studies the \emph{correctness} or \emph{validity} of the
relations of reason we construct.

However, studying that aspect of logic alone is incomplete. To
illustrate the problem, consider the following syllogism:

\begin{quote}
All animals have wings.\\
All Frogs are animals.\\
Therefore, all frogs have wings.
\end{quote}

\noindent The syllogism is \emph{formally} perfectly correct. As you
will learn in the logic of the argument, the middle term, \emph{animal},
is universal at least once, which is sufficient to join the major and
minor premise. In short, it is formally a perfect syllogism in BARBARA.
However, the conclusion is evidently false. The reason, of course, is
that the major premise (``All animals have wings'') is false. (In fact,
only \emph{some} animals have wings, but not all.) Thus, as you can see,
besides formal correctness, there is an additional condition to reaching
a true conclusion: both premises must themselves be true. The conditions
under which we can guarantee that \emph{formally} correct reasoning is
not merely valid but leads to \emph{true conclusions} are called the
\emph{material conditions} for truth, and their study is called
\emph{material logic} or \emph{major logic}.

As my example makes clear, it is easy to apply the division between
material and formal logic to the logic of the argument. Does it apply
equally well to the other parts of logic? For the logic of the concept,
it does in a way. concepts are called \emph{universals},\footnote{Traditionally,
  it is said that etymologically the word ``universal'' comes from the
  Latin phrase \emph{unum versus alia}, ``one toward (many) others.''}
because through one concept, we are able to know many concrete
individuals. We can, therefore two aspects of this universality: there
is the \emph{manner} in which the concept is attributed universally to
its subject, which we call \emph{predicability}\footnote{You will learn
  in the course that there are five such ways, or predicables: the
  genus, the species, the difference, the property, and the (predicable)
  accident.}, and there are the \emph{kinds of objects} (or
\emph{contents}) that the concept can refer to.\footnote{You will learn
  in the course that these types of objects include the transcendentals
  (the properties of all entities), the categories, and the
  post-predicaments.} In a sense, the latter aspect constitutes the
``material'' conditions for the truth of the concept, and thus by
analogy, the former constitutes the ``formal'' conditions, because it
leaves aside the question of \emph{what} the concept is referring to.
Thus, with the logic of the concept, the division between formal and
material logic is somewhat justifiable---even though the criterion for
dividing is rather different.

It is much more problematic to apply this division to the logic of the
enunciation. As we have already discussed, when we attribute a predicate
to a subject (or separate it from that subject), we can do so either
mediately---that is, with the help of other enunciations that form our
premises---or immediately---that is, by understanding it directly either
from the first principles of reason or from experience. Of course, we
obtain mediated knowledge through argument, and thus, and we know that
the logic of the argument is rightly divided between formal and material
logic. What is left is \emph{immediate} knowledge of a real object:
however, by its very nature, immediate knowledge is not subject to
questions about conditions for truth: it is already \emph{known} to be
true.

\chapter*{Conclusion}\label{conclusion}
\addcontentsline{toc}{chapter}{Conclusion}

As good logicians, we wished to ask ourselves, with precision and rigor,
``What precisely is logic?'' We are now in a much better position to
answer that question, having explored the place of logic among the other
arts and sciences, and even explored exactly how it perfects the
intellect as an intellectual \emph{habitus}. You should now be in a
position to explain the definition of logic, and especially what its
object of study is: the works of reason as regards the relations that we
construct between second-order intentions, in the light the first
principles of reason. You should now be better prepared to explore that
object, in the three great divisions of this science: the logic of the
concept, the logic of the enunciation, and the logic of the argument.

\backmatter
\end{document}
